\section{Source atlas}
\label{app:atlas}

\AtlasRecords{} source records: what each file is, what is located inside
it, and the scope within which it may be used.

Two of those fields --- \textit{Located} and \textit{Scope} --- are
curated judgement, carried over from a companion edition of this
program (\S\ref{sec:lineage}). Everything else on each row is
re-derived when this document is built: the file is looked for, its
bytes are hashed against the recorded digest, and its graph identifiers
are looked up in the catalogue. A record therefore carries a state, and
the states mean different things:

\begin{description}[leftmargin=1.2em,style=nextline]
\item[verified] The file is in this checkout and its bytes are the ones
  the note was written against. \AtlasVerified{} of \AtlasRecords{}.
\item[changed since curated] The file is here and its bytes have moved
  on. This is expected for the maintained ledgers and the generated
  catalogues, which change whenever research lands; it is a signal to
  re-read the note, not evidence of an error.
\item[outside this checkout] A real source in the wider workspace
  archive or in a sibling worktree. This build cannot hash it. These
  are among the oldest and most load-bearing material in the program
  and their absence here is a property of the repository boundary, not
  of the source.
\item[not on this branch] A repository path this branch does not carry
  --- in practice, a source still on an open pull request. The curation
  is ahead of \texttt{main} rather than wrong.
\end{description}

That split is the reason this appendix is generated rather than
transcribed. A hand-maintained atlas states a digest once and goes on
stating it; this one re-checks on every build, so the drift is visible
in the document instead of waiting to mislead someone.

% Generated by paper/master/generate_sources.py.  Do not edit.
\noindent 123 source records. Each names a local file, what is located in it, and the scope within which it may be used. The first two fields are curated judgement; the rest is re-derived at generation time --- the file is hashed and its identifiers are looked up in the catalogue, so a moved file or a renamed claim appears below as a marked row instead of as a sentence that quietly stopped being true.\par\medskip
\begingroup\sloppy
\noindent\textbf{\texttt{adr0024}} \hfill {\footnotesize verified}\par
{\small ADR 0024: The corner cluster from a third implementation, and the historical ledger that was here all along}\par
{\footnotesize\ttfamily C:/\allowbreak{}WORKHOUSE/\allowbreak{}REPO/\allowbreak{}docs/\allowbreak{}decisions/\allowbreak{}0024-\allowbreak{}the-\allowbreak{}corner-\allowbreak{}cluster-\allowbreak{}from-\allowbreak{}a-\allowbreak{}third-\allowbreak{}implementation-\allowbreak{}and-\allowbreak{}the-\allowbreak{}ledger-\allowbreak{}that-\allowbreak{}was-\allowbreak{}here.md}\par
{\footnotesize \textit{Located.} What was found 1-4; Decision C2}\par
{\footnotesize \textit{Scope.} Resolved C\_\allowbreak{}shp and the sixteen missing adjacent-face cube orderings}\par
{\footnotesize SHA-256 \texttt{d097de0151c13ff6} \quad graph: \texttt{ADR:0024}}\par
\smallskip
\noindent\textbf{\texttt{adr0029}} \hfill {\footnotesize verified}\par
{\small ADR 0029: The degree bound is proved by removing it: the third engine over Q(N), and the corner and fan by channel}\par
{\footnotesize\ttfamily C:/\allowbreak{}WORKHOUSE/\allowbreak{}REPO/\allowbreak{}docs/\allowbreak{}decisions/\allowbreak{}0029-\allowbreak{}the-\allowbreak{}degree-\allowbreak{}bound-\allowbreak{}is-\allowbreak{}proved-\allowbreak{}by-\allowbreak{}removing-\allowbreak{}it.md}\par
{\footnotesize \textit{Located.} What was found 1-5; What is not claimed}\par
{\footnotesize \textit{Scope.} Exact rational closed cumulants, N>=5 specialization argument, Lean identity on literal forms, retained model-space premises}\par
{\footnotesize SHA-256 \texttt{d1032c4477bc35c6} \quad graph: \texttt{ADR:0029}}\par
\smallskip
\noindent\textbf{\texttt{adr0046}} \hfill {\footnotesize verified}\par
{\small ADR 0046: The fourth-order band is N\textasciicircum{}-7 because two orders cancel in every cluster, and the band is rank-uniform in tau through fourth order}\par
{\footnotesize\ttfamily C:/\allowbreak{}WORKHOUSE/\allowbreak{}REPO/\allowbreak{}docs/\allowbreak{}decisions/\allowbreak{}0046-\allowbreak{}the-\allowbreak{}fourth-\allowbreak{}order-\allowbreak{}band-\allowbreak{}is-\allowbreak{}planar-\allowbreak{}suppressed-\allowbreak{}two-\allowbreak{}orders-\allowbreak{}below-\allowbreak{}its-\allowbreak{}channels.md}\par
{\footnotesize \textit{Located.} Decision facts 1-3 and conjecture 4}\par
{\footnotesize \textit{Scope.} Sixteen recorded cluster/sector pairs and 1772 channels, fixed order; general-k and G16 overlap open}\par
{\footnotesize SHA-256 \texttt{5f3a1e07b42059f7} \quad graph: \texttt{ADR:0046}}\par
\smallskip
\noindent\textbf{\texttt{adr0047}} \hfill {\footnotesize verified}\par
{\small The odd orders of the strong-coupling band are determinant families}\par
{\footnotesize\ttfamily C:/\allowbreak{}WORKHOUSE/\allowbreak{}REPO/\allowbreak{}paper/\allowbreak{}research\_\allowbreak{}notes/\allowbreak{}ODD\_\allowbreak{}ORDER\_\allowbreak{}CENTRE\_\allowbreak{}PARITY\_\allowbreak{}20260911.md}\par
{\footnotesize \textit{Located.} Sections 2-3 proofs; Section 6 finite-cluster/even-periodic-extent scope}\par
{\footnotesize \textit{Scope.} All-order centre parity coefficient theorem, finite clusters of Z3 or even torus extents; exact first odd vertex}\par
{\footnotesize SHA-256 \texttt{7a25754b96c10a00} \quad graph: \texttt{CITE:\allowbreak{}ODD\_\allowbreak{}ORDER\_\allowbreak{}CENTRE\_\allowbreak{}PARITY}}\par
\smallskip
\noindent\textbf{\texttt{adr0048}} \hfill {\footnotesize not found}\par
{\small The single-face planar grading law: orders eight and ten, and the irrep channels that carry the cancellation}\par
{\footnotesize\ttfamily C:/\allowbreak{}WORKHOUSE/\allowbreak{}REPO/\allowbreak{}paper/\allowbreak{}research\_\allowbreak{}notes/\allowbreak{}PLANAR\_\allowbreak{}GRADING\_\allowbreak{}SINGLE\_\allowbreak{}FACE\_\allowbreak{}20260911.md}\par
{\footnotesize \textit{Located.} Object; Results 1-3; What this does and does not establish}\par
{\footnotesize \textit{Scope.} One-face rank-generic rational coefficients through m10, locally developed open PR164 at pinned snapshot. General-m and cluster law not proved. Order12 observed outside certification}\par
{\footnotesize SHA-256 \texttt{cecfb472a318698f} \quad graph: \texttt{CITE:\allowbreak{}PLANAR\_\allowbreak{}GRADING\_\allowbreak{}SINGLE\_\allowbreak{}FACE}}\par
\smallskip
\noindent\textbf{\texttt{anc-archive-manifest}} \hfill {\footnotesize verified}\par
{\small Historical-archive original-source manifest}\par
{\footnotesize\ttfamily C:/\allowbreak{}WORKHOUSE/\allowbreak{}REPO/\allowbreak{}runs/\allowbreak{}archive\_\allowbreak{}derivations\_\allowbreak{}2026-\allowbreak{}09-\allowbreak{}11/\allowbreak{}source-\allowbreak{}manifest.json}\par
{\footnotesize \textit{Located.} Complete manifest}\par
{\footnotesize \textit{Scope.} Maps original Appendix E, Parts7-8, Appendices G-H and source-tilt material to retained copies and hashes}\par
{\footnotesize SHA-256 \texttt{01f408e56c5cbef0}}\par
\smallskip
\noindent\textbf{\texttt{anc-earlier}} \hfill {\footnotesize verified}\par
{\small Recovered earlier proofs and their exact application boundaries}\par
{\footnotesize\ttfamily C:/\allowbreak{}WORKHOUSE/\allowbreak{}REPO/\allowbreak{}docs/\allowbreak{}derivations/\allowbreak{}earlier-\allowbreak{}proof-\allowbreak{}recovery.md}\par
{\footnotesize \textit{Located.} E1, lines 14-37; E3, lines 54-74; E5a-E5d, lines 103-167}\par
{\footnotesize \textit{Scope.} Analytic recovered theorems with explicit application hypotheses; original arguments by Alexander Smith}\par
{\footnotesize SHA-256 \texttt{4b5b067ee7697cf5} \quad graph: \texttt{CITE:\allowbreak{}EARLIER\_\allowbreak{}PROOF\_\allowbreak{}RECOVERY}}\par
\smallskip
\noindent\textbf{\texttt{anc-earlier-manifest}} \hfill {\footnotesize verified}\par
{\small Earlier-proof recovery original-source manifest}\par
{\footnotesize\ttfamily C:/\allowbreak{}WORKHOUSE/\allowbreak{}REPO/\allowbreak{}runs/\allowbreak{}earlier\_\allowbreak{}proof\_\allowbreak{}recovery\_\allowbreak{}2026-\allowbreak{}09-\allowbreak{}11/\allowbreak{}source\_\allowbreak{}manifest.json}\par
{\footnotesize \textit{Located.} Complete manifest}\par
{\footnotesize \textit{Scope.} Maps earlier original sources and retained copies to hashes}\par
{\footnotesize SHA-256 \texttt{cf5fcb2eaa6f6bdb}}\par
\smallskip
\noindent\textbf{\texttt{anc-functional}} \hfill {\footnotesize verified}\par
{\small Smooth defects, Lyapunov patching and concentration from the historical archive}\par
{\footnotesize\ttfamily C:/\allowbreak{}WORKHOUSE/\allowbreak{}REPO/\allowbreak{}docs/\allowbreak{}derivations/\allowbreak{}archive-\allowbreak{}functional-\allowbreak{}inequalities.md}\par
{\footnotesize \textit{Located.} F2, lines 23-36; F8-F9, lines 115-143; F17, lines 271-310}\par
{\footnotesize \textit{Scope.} Analytic conditional functional estimates and exact unit-S3 SU2 signed diffusion identity; physical transfer matching remains separate}\par
{\footnotesize SHA-256 \texttt{048fbf01cdc55ee4} \quad graph: \texttt{CITE:\allowbreak{}ARCHIVE\_\allowbreak{}FUNCTIONAL}}\par
\smallskip
\noindent\textbf{\texttt{anc-gpu}} \hfill {\footnotesize outside this checkout}\par
{\small GPU TESTS: reproduced calculations and theory-graph audit}\par
{\footnotesize\ttfamily C:/\allowbreak{}WORKHOUSE/\allowbreak{}worktrees/\allowbreak{}gpu-\allowbreak{}tests-\allowbreak{}audit-\allowbreak{}20260912/\allowbreak{}docs/\allowbreak{}derivations/\allowbreak{}gpu-\allowbreak{}tests-\allowbreak{}audit.md}\par
{\footnotesize \textit{Located.} A1, lines 280-290; A4, lines 310-320}\par
{\footnotesize \textit{Scope.} Local unpublished audit: pointwise specified-chart SU3 identity Hessian, all-size specified scalar-chain Schur floor}\par
{\footnotesize SHA-256 \texttt{6e95c24f725896aa} \quad \textbf{workspace path under worktrees/}}\par
\smallskip
\noindent\textbf{\texttt{anc-gpu-maxwell}} \hfill {\footnotesize outside this checkout}\par
{\small GPU Tests 26-30: dated continuation of the source audit}\par
{\footnotesize\ttfamily C:/\allowbreak{}WORKHOUSE/\allowbreak{}worktrees/\allowbreak{}gpu-\allowbreak{}tests-\allowbreak{}audit-\allowbreak{}20260912/\allowbreak{}docs/\allowbreak{}derivations/\allowbreak{}gpu-\allowbreak{}tests-\allowbreak{}followup-\allowbreak{}26-\allowbreak{}30.md}\par
{\footnotesize \textit{Located.} T29}\par
{\footnotesize \textit{Scope.} T2 finite correct-operator curl-curl Maxwell sweep on 3\textasciicircum{}4 torus, all pairs, eight parameter pairs; rational inverse and Arb comparison}\par
{\footnotesize SHA-256 \texttt{b34105f1fd5bf099} \quad \textbf{workspace path under worktrees/}}\par
\smallskip
\noindent\textbf{\texttt{anc-gpu-record}} \hfill {\footnotesize outside this checkout}\par
{\small GPU Tests audit: completed combined-tree validation}\par
{\footnotesize\ttfamily C:/\allowbreak{}WORKHOUSE/\allowbreak{}worktrees/\allowbreak{}gpu-\allowbreak{}tests-\allowbreak{}audit-\allowbreak{}20260912/\allowbreak{}graph-\allowbreak{}tasks/\allowbreak{}2026-\allowbreak{}09-\allowbreak{}12-\allowbreak{}gpu-\allowbreak{}tests-\allowbreak{}audit.md}\par
{\footnotesize \textit{Located.} Completed scientific registration and final validation; Completed combined-tree validation}\par
{\footnotesize \textit{Scope.} All 35 source programs reviewed and replayed; 40 statements, 22 new controls; local combined 704/704 checks and 2144 passed, 4 skipped. Publication not performed}\par
{\footnotesize SHA-256 \texttt{2a0c490c59867b47} \quad \textbf{workspace path under worktrees/}}\par
\smallskip
\noindent\textbf{\texttt{anc-prior-final}} \hfill {\footnotesize verified}\par
{\small Volume-uniform electric-shell isolation, exact shared-link hopping, and a homological carrier in strong-coupling SU(N) Hamiltonian lattice gauge theory}\par
{\footnotesize\ttfamily C:/\allowbreak{}Users/\allowbreak{}Alex/\allowbreak{}Desktop/\allowbreak{}FINAL PAPERS/\allowbreak{}WORKHOUSE\_\allowbreak{}PUBLICATION\_\allowbreak{}EDITION\_\allowbreak{}FINAL\_\allowbreak{}20260830/\allowbreak{}workhouse\_\allowbreak{}publication\_\allowbreak{}edition\_\allowbreak{}final\_\allowbreak{}20260830.tex}\par
{\footnotesize \textit{Located.} Title/author/date, lines 45-49; Abstract, lines 52-134; Introduction claims, lines 179-313; Fourth-order boundary, line 2273 onward; Fixed-spacing band, line 3008 onward; Scope, line 3012 onward}\par
{\footnotesize \textit{Scope.} User-authored earlier manuscript; retained as historical lineage, not current theory-graph authority for superseded O4-open claims}\par
{\footnotesize SHA-256 \texttt{7d9022c8c4d49d95}}\par
\smallskip
\noindent\textbf{\texttt{anc-transport}} \hfill {\footnotesize verified}\par
{\small Source tilts, reflected pushforwards and finite-range inverse transport}\par
{\footnotesize\ttfamily C:/\allowbreak{}WORKHOUSE/\allowbreak{}REPO/\allowbreak{}docs/\allowbreak{}derivations/\allowbreak{}archive-\allowbreak{}source-\allowbreak{}transport.md}\par
{\footnotesize \textit{Located.} S1-S4, lines 10-56; S7-S8, lines 95-121}\par
{\footnotesize \textit{Scope.} Bounded source identities, normalized domination and reflection-positive permanence under positive-time pullback and convergence hypotheses}\par
{\footnotesize SHA-256 \texttt{4a9d953d888e1b0e} \quad graph: \texttt{CITE:\allowbreak{}ARCHIVE\_\allowbreak{}SOURCE\_\allowbreak{}TRANSPORT}}\par
\smallskip
\noindent\textbf{\texttt{anisotropy}} \hfill {\footnotesize verified}\par
{\small Geometry, source transport, and residual certificates in the theory graph}\par
{\footnotesize\ttfamily C:/\allowbreak{}WORKHOUSE/\allowbreak{}REPO/\allowbreak{}docs/\allowbreak{}derivations/\allowbreak{}theory-\allowbreak{}geometry-\allowbreak{}and-\allowbreak{}riccati-\allowbreak{}bridges.md}\par
{\footnotesize \textit{Located.} TG1 sharp three-coordinate bound; TG2 zero extension; TG3 exact Hodge-support secular cubic}\par
{\footnotesize \textit{Scope.} Analytic sharp K3 bound and exact 3x3 secular identity, with source-specific machine coverage}\par
{\footnotesize SHA-256 \texttt{7216dd8c2f263782} \quad graph: \texttt{CITE:\allowbreak{}THEORY\_\allowbreak{}GEOMETRY\_\allowbreak{}RICCATI\_\allowbreak{}BRIDGES}}\par
\smallskip
\noindent\textbf{\texttt{contAntipodal}} \hfill {\footnotesize verified}\par
{\small W6 antipodal magnetic geometry and the remaining conditional score}\par
{\footnotesize\ttfamily C:/\allowbreak{}WORKHOUSE/\allowbreak{}REPO/\allowbreak{}docs/\allowbreak{}derivations/\allowbreak{}w6-\allowbreak{}antipodal-\allowbreak{}magnetic-\allowbreak{}geometry.md}\par
{\footnotesize \textit{Located.} A1-A7}\par
{\footnotesize SHA-256 \texttt{5b3e0b2312f11ee8} \quad graph: \texttt{CITE:\allowbreak{}W6\_\allowbreak{}ANTIPODAL\_\allowbreak{}MAGNETIC\_\allowbreak{}GEOMETRY}}\par
\smallskip
\noindent\textbf{\texttt{contBudgets}} \hfill {\footnotesize verified}\par
{\small Corrected fourth-order coercivity and sharp multichannel spectral budgets}\par
{\footnotesize\ttfamily C:/\allowbreak{}WORKHOUSE/\allowbreak{}REPO/\allowbreak{}docs/\allowbreak{}derivations/\allowbreak{}review-\allowbreak{}derived-\allowbreak{}spectral-\allowbreak{}budgets.md}\par
{\footnotesize \textit{Located.} Q1-Q2 corrected coercivity; M1-M3 multichannel spectral budgets}\par
{\footnotesize SHA-256 \texttt{96037da9bfc02943} \quad graph: \texttt{CITE:\allowbreak{}REVIEW\_\allowbreak{}SPECTRAL\_\allowbreak{}BUDGETS}}\par
\smallskip
\noindent\textbf{\texttt{contComplement}} \hfill {\footnotesize verified}\par
{\small Scoped results and actual complement control from the M10 proof repair}\par
{\footnotesize\ttfamily C:/\allowbreak{}WORKHOUSE/\allowbreak{}REPO/\allowbreak{}paper/\allowbreak{}research\_\allowbreak{}notes/\allowbreak{}W6\_\allowbreak{}M10\_\allowbreak{}REPAIR\_\allowbreak{}RESULTS\_\allowbreak{}20260911.md}\par
{\footnotesize \textit{Located.} Actual centered conditional complement theorem}\par
{\footnotesize SHA-256 \texttt{d57148c48497b175} \quad graph: \texttt{CITE:\allowbreak{}W6\_\allowbreak{}M10\_\allowbreak{}REPAIR\_\allowbreak{}RESULTS}}\par
\smallskip
\noindent\textbf{\texttt{contCurrents}} \hfill {\footnotesize verified}\par
{\small Complete square-source currents and spectral totality}\par
{\footnotesize\ttfamily C:/\allowbreak{}WORKHOUSE/\allowbreak{}REPO/\allowbreak{}docs/\allowbreak{}derivations/\allowbreak{}source-\allowbreak{}currents-\allowbreak{}and-\allowbreak{}spectral-\allowbreak{}totality.md}\par
{\footnotesize \textit{Located.} SCB0-SCB3}\par
{\footnotesize SHA-256 \texttt{860c81a3fb665d6b} \quad graph: \texttt{CITE:\allowbreak{}SOURCE\_\allowbreak{}CURRENT\_\allowbreak{}SPECTRAL\_\allowbreak{}BRIDGES}}\par
\smallskip
\noindent\textbf{\texttt{contG18}} \hfill {\footnotesize verified}\par
{\small The complete infinite-volume physical Wilson band and literal-source frame}\par
{\footnotesize\ttfamily C:/\allowbreak{}WORKHOUSE/\allowbreak{}REPO/\allowbreak{}docs/\allowbreak{}derivations/\allowbreak{}wilson-\allowbreak{}shell-\allowbreak{}inputs/\allowbreak{}G18\_\allowbreak{}WILSON\_\allowbreak{}INFINITE\_\allowbreak{}VOLUME\_\allowbreak{}PHYSICAL\_\allowbreak{}BAND\_\allowbreak{}20260905.md}\par
{\footnotesize \textit{Located.} 1, 4-9}\par
{\footnotesize SHA-256 \texttt{966e464c641e96aa} \quad graph: \texttt{CITE:\allowbreak{}WILSON\_\allowbreak{}BAND\_\allowbreak{}INPUT}}\par
\smallskip
\noindent\textbf{\texttt{contJets}} \hfill {\footnotesize verified}\par
{\small W6: raw derivative scaling, the remaining transport estimate, and the physical budget}\par
{\footnotesize\ttfamily C:/\allowbreak{}WORKHOUSE/\allowbreak{}REPO/\allowbreak{}docs/\allowbreak{}derivations/\allowbreak{}w6-\allowbreak{}ground-\allowbreak{}jets-\allowbreak{}and-\allowbreak{}transport-\allowbreak{}budget.md}\par
{\footnotesize \textit{Located.} R8a-R15; section 6}\par
{\footnotesize SHA-256 \texttt{a1638b09ba941b0c} \quad graph: \texttt{CITE:\allowbreak{}W6\_\allowbreak{}GROUND\_\allowbreak{}JETS\_\allowbreak{}TRANSPORT\_\allowbreak{}BUDGET}}\par
\smallskip
\noindent\textbf{\texttt{contKernels}} \hfill {\footnotesize verified}\par
{\small A kernel criterion for the G19 moving-time route}\par
{\footnotesize\ttfamily C:/\allowbreak{}WORKHOUSE/\allowbreak{}REPO/\allowbreak{}docs/\allowbreak{}derivations/\allowbreak{}os-\allowbreak{}kernel-\allowbreak{}moving-\allowbreak{}time-\allowbreak{}gap.md}\par
{\footnotesize \textit{Located.} K1-K5}\par
{\footnotesize SHA-256 \texttt{d7088d36f67a1f8c} \quad graph: \texttt{CITE:\allowbreak{}OS\_\allowbreak{}KERNEL\_\allowbreak{}GAP}}\par
\smallskip
\noindent\textbf{\texttt{contM10}} \hfill {\footnotesize verified}\par
{\small W6 synchronized M10: repaired criteria and remaining actual-ground estimate}\par
{\footnotesize\ttfamily C:/\allowbreak{}WORKHOUSE/\allowbreak{}REPO/\allowbreak{}docs/\allowbreak{}derivations/\allowbreak{}w6-\allowbreak{}synchronized-\allowbreak{}m10-\allowbreak{}domination.md}\par
{\footnotesize \textit{Located.} 1-7}\par
{\footnotesize SHA-256 \texttt{e81031e9f7d13467} \quad graph: \texttt{CITE:\allowbreak{}W6\_\allowbreak{}SYNCHRONIZED\_\allowbreak{}M10\_\allowbreak{}DOMINATION}}\par
\smallskip
\noindent\textbf{\texttt{contMoving}} \hfill {\footnotesize verified}\par
{\small Moving-time spectral exclusion with approximate sources}\par
{\footnotesize\ttfamily C:/\allowbreak{}WORKHOUSE/\allowbreak{}REPO/\allowbreak{}docs/\allowbreak{}derivations/\allowbreak{}moving-\allowbreak{}time-\allowbreak{}spectral-\allowbreak{}gap.md}\par
{\footnotesize \textit{Located.} MT1-MT5}\par
{\footnotesize SHA-256 \texttt{d0a20c25e023bff7} \quad graph: \texttt{CITE:\allowbreak{}MOVING\_\allowbreak{}TIME\_\allowbreak{}SPECTRAL\_\allowbreak{}GAP}}\par
\smallskip
\noindent\textbf{\texttt{contR10}} \hfill {\footnotesize verified}\par
{\small W6: repaired source-energy identities and the remaining R10 estimates}\par
{\footnotesize\ttfamily C:/\allowbreak{}WORKHOUSE/\allowbreak{}REPO/\allowbreak{}docs/\allowbreak{}derivations/\allowbreak{}w6-\allowbreak{}source-\allowbreak{}energy-\allowbreak{}jets-\allowbreak{}r10.md}\par
{\footnotesize \textit{Located.} C1-C6}\par
{\footnotesize SHA-256 \texttt{7cb2ad710a35e1dd} \quad graph: \texttt{CITE:\allowbreak{}W6\_\allowbreak{}SOURCE\_\allowbreak{}ENERGY\_\allowbreak{}JETS\_\allowbreak{}R10}}\par
\smallskip
\noindent\textbf{\texttt{contSC17}} \hfill {\footnotesize verified}\par
{\small SC17: connected spatial closure and the signed background continuation}\par
{\footnotesize\ttfamily C:/\allowbreak{}WORKHOUSE/\allowbreak{}REPO/\allowbreak{}docs/\allowbreak{}derivations/\allowbreak{}wilson-\allowbreak{}sc17-\allowbreak{}spatial-\allowbreak{}closure.md}\par
{\footnotesize \textit{Located.} Introduction; optimized interval and endpoint}\par
{\footnotesize SHA-256 \texttt{d02e906ba0d52d20} \quad graph: \texttt{CITE:\allowbreak{}WILSON\_\allowbreak{}SC17}}\par
\smallskip
\noindent\textbf{\texttt{contSC17Time}} \hfill {\footnotesize verified}\par
{\small SC17 consequence: actual infinite-volume physical Euclidean-time semigroup}\par
{\footnotesize\ttfamily C:/\allowbreak{}WORKHOUSE/\allowbreak{}REPO/\allowbreak{}docs/\allowbreak{}derivations/\allowbreak{}wilson-\allowbreak{}sc17-\allowbreak{}physical-\allowbreak{}time-\allowbreak{}limit.md}\par
{\footnotesize \textit{Located.} P1-P10}\par
{\footnotesize SHA-256 \texttt{f842a6b9f826a997} \quad graph: \texttt{CITE:\allowbreak{}WILSON\_\allowbreak{}SC17\_\allowbreak{}TIME}}\par
\smallskip
\noindent\textbf{\texttt{contScale}} \hfill {\footnotesize verified}\par
{\small Wilson spatial passage: complete Schur excess and scale budgets}\par
{\footnotesize\ttfamily C:/\allowbreak{}WORKHOUSE/\allowbreak{}REPO/\allowbreak{}docs/\allowbreak{}derivations/\allowbreak{}wilson-\allowbreak{}spatial-\allowbreak{}schur-\allowbreak{}excess.md}\par
{\footnotesize \textit{Located.} Section 6 SP20-SP24; Section 7 SP25 and the normalized-coarse-energy requirement}\par
{\footnotesize \textit{Scope.} Conditional accumulation of complete physical gaps, source frames and specified-observable amplitudes. Actual Wilson matching and continuum correlation limits remain hypotheses}\par
{\footnotesize SHA-256 \texttt{5ddb7f812a9e6f53} \quad graph: \texttt{CITE:\allowbreak{}WILSON\_\allowbreak{}SPATIAL}}\par
\smallskip
\noindent\textbf{\texttt{contScore}} \hfill {\footnotesize verified}\par
{\small W6: the complete source generator in the ground-state frame and the score Poisson equation}\par
{\footnotesize\ttfamily C:/\allowbreak{}WORKHOUSE/\allowbreak{}REPO/\allowbreak{}docs/\allowbreak{}derivations/\allowbreak{}w6-\allowbreak{}source-\allowbreak{}generator-\allowbreak{}score-\allowbreak{}frame.md}\par
{\footnotesize \textit{Located.} SF1-SF9}\par
{\footnotesize SHA-256 \texttt{20b070cb78311187} \quad graph: \texttt{CITE:\allowbreak{}W6\_\allowbreak{}SOURCE\_\allowbreak{}GENERATOR\_\allowbreak{}SCORE\_\allowbreak{}FRAME}}\par
\smallskip
\noindent\textbf{\texttt{contScoreTail}} \hfill {\footnotesize verified}\par
{\small Actual compact-square conditional-score tail bound}\par
{\footnotesize\ttfamily C:/\allowbreak{}WORKHOUSE/\allowbreak{}REPO/\allowbreak{}docs/\allowbreak{}derivations/\allowbreak{}w6-\allowbreak{}conditional-\allowbreak{}score-\allowbreak{}tail-\allowbreak{}control.md}\par
{\footnotesize \textit{Located.} Section 6, M.2-M.5, M5-M15: source-potential moment and conditional M10-to-weighted-score/Hardy implication}\par
{\footnotesize \textit{Scope.} Actual source-potential estimate and a conditional implication retaining the tube-moment, relative-amplitude and normal/angular assumptions}\par
{\footnotesize SHA-256 \texttt{d04ae01d29d33cc5} \quad graph: \texttt{CITE:\allowbreak{}W6\_\allowbreak{}CONDITIONAL\_\allowbreak{}SCORE\_\allowbreak{}TAIL\_\allowbreak{}CONTROL}}\par
\smallskip
\noindent\textbf{\texttt{contSelected}} \hfill {\footnotesize verified}\par
{\small The full-space Gaussian test, its repair, and the precise remaining estimate}\par
{\footnotesize\ttfamily C:/\allowbreak{}WORKHOUSE/\allowbreak{}REPO/\allowbreak{}docs/\allowbreak{}derivations/\allowbreak{}wilson-\allowbreak{}selected-\allowbreak{}inverse-\allowbreak{}wall.md}\par
{\footnotesize \textit{Located.} 3-5}\par
{\footnotesize SHA-256 \texttt{c0473aca76790ef7} \quad graph: \texttt{CITE:\allowbreak{}WILSON\_\allowbreak{}SELECTED}}\par
\smallskip
\noindent\textbf{\texttt{evenband}} \hfill {\footnotesize verified}\par
{\small The charge-even band, exactly}\par
{\footnotesize\ttfamily C:/\allowbreak{}WORKHOUSE/\allowbreak{}REPO/\allowbreak{}src/\allowbreak{}workhouse/\allowbreak{}invariants/\allowbreak{}charge\_\allowbreak{}even.py}\par
{\footnotesize \textit{Located.} Checks: characteristic polynomial (lines 105-123), floor planes (lines 224-258), one assembly formula (lines 299-328), and Gamma curvature}\par
{\footnotesize \textit{Scope.} Exact unsigned-incidence Bloch cubic, floor planes and multiplicities, and SU3 charge-even band assembly through order three. Do not use unique-floor prose; dedicated floor check gives three planes}\par
{\footnotesize SHA-256 \texttt{a6591f321b6cd2c6}}\par
\smallskip
\noindent\textbf{\texttt{evenband-support}} \hfill {\footnotesize verified}\par
{\small Charge-even incidence implementation}\par
{\footnotesize\ttfamily C:/\allowbreak{}WORKHOUSE/\allowbreak{}REPO/\allowbreak{}src/\allowbreak{}workhouse/\allowbreak{}even\_\allowbreak{}sector.py}\par
{\footnotesize \textit{Located.} See source record}\par
{\footnotesize SHA-256 \texttt{48c7d96fc26cc3c2}}\par
\smallskip
\noindent\textbf{\texttt{fourth20260909}} \hfill {\footnotesize verified}\par
{\small The fourth order at every rank: closed forms, a Hodge decomposition, and the mechanism of the tier collapse in strong-coupling SU(N) Hamiltonian lattice gauge theory}\par
{\footnotesize\ttfamily C:/\allowbreak{}WORKHOUSE/\allowbreak{}REPO/\allowbreak{}paper/\allowbreak{}fourth\_\allowbreak{}order\_\allowbreak{}all\_\allowbreak{}rank\_\allowbreak{}2026-\allowbreak{}09-\allowbreak{}09.tex}\par
{\footnotesize \textit{Located.} Introduction: object, premises, second order, shape ansatz; Hodge and support mechanism; all-rank assembly}\par
{\footnotesize \textit{Scope.} Fixed-spacing projected charge-odd shell; four recorded engine/support premises; all-rank specialization argument N>=5, exact comparisons N=3,4}\par
{\footnotesize SHA-256 \texttt{0681a367d8d65e86}}\par
\smallskip
\noindent\textbf{\texttt{fourthAnalyticTemporal}} \hfill {\footnotesize verified}\par
{\small Fixed-observable Wilson expansion, complete shell, and carrier transport}\par
{\footnotesize\ttfamily C:/\allowbreak{}WORKHOUSE/\allowbreak{}REPO/\allowbreak{}docs/\allowbreak{}derivations/\allowbreak{}wilson-\allowbreak{}marked-\allowbreak{}shell-\allowbreak{}transport.md}\par
{\footnotesize \textit{Located.} Section 1 setup and S1; Sections 2-4 S7-S15 complex operator and frame construction; Sections 5-7 S16-S22 weighted majorant, summed temporal matching, relative carrier and source amplitude}\par
{\footnotesize \textit{Scope.} Fixed-spatial-spacing SU(3), admitted temporal Wilson meshes and common small-magnetic disc. Weighted analytic majorant and conditional temporal Hamiltonian/Wilson matching retain the source window and Hamiltonian relative-gap hypotheses; no spatial continuum or isotropic identification is inferred}\par
{\footnotesize SHA-256 \texttt{a887ed9de154cec0} \quad graph: \texttt{CITE:\allowbreak{}WILSON\_\allowbreak{}SHELL}}\par
\smallskip
\noindent\textbf{\texttt{fourthAnalyticThermo}} \hfill {\footnotesize verified}\par
{\small SC17: thermodynamic vacuum, closed ground-law form and plaquette overlap}\par
{\footnotesize\ttfamily C:/\allowbreak{}WORKHOUSE/\allowbreak{}REPO/\allowbreak{}docs/\allowbreak{}derivations/\allowbreak{}wilson-\allowbreak{}sc17-\allowbreak{}thermodynamic-\allowbreak{}limit.md}\par
{\footnotesize \textit{Located.} Part I T1-T13 and uniform score limits; Part II IF10-IF13, with IF4-IF9 physical form and gap}\par
{\footnotesize \textit{Scope.} Actual fixed-spacing SU(2) normalized ground response and clustering in the SC17 interior, free/periodic thermodynamic limit including the endpoint, explicit nonzero physical plaquette interval weight}\par
{\footnotesize SHA-256 \texttt{6d72876a35a86676} \quad graph: \texttt{CITE:\allowbreak{}WILSON\_\allowbreak{}SC17\_\allowbreak{}LIMIT}}\par
\smallskip
\noindent\textbf{\texttt{fourthCapChecks}} \hfill {\footnotesize verified}\par
{\small Exact isotropic pentagonal cap-band controls}\par
{\footnotesize\ttfamily C:/\allowbreak{}WORKHOUSE/\allowbreak{}REPO/\allowbreak{}src/\allowbreak{}workhouse/\allowbreak{}invariants/\allowbreak{}pentagonal.py}\par
{\footnotesize \textit{Located.} Checks for electric energies, endpoint difference, D5 symmetry, Hermitian symbol, bandwidth and target-blind cold record}\par
{\footnotesize \textit{Scope.} Exact assembly identities and retained-source controls; fifth-order physical histories are not re-executed by these checks}\par
{\footnotesize SHA-256 \texttt{6d2994bdbd721afc}}\par
\smallskip
\noindent\textbf{\texttt{fourthCapMaster}} \hfill {\footnotesize verified}\par
{\small Unified master theory, version 4.3}\par
{\footnotesize\ttfamily C:/\allowbreak{}WORKHOUSE/\allowbreak{}REPO/\allowbreak{}theory/\allowbreak{}MASTER\_\allowbreak{}THEORY\_\allowbreak{}UNIFIED\_\allowbreak{}2026-\allowbreak{}08-\allowbreak{}20\_\allowbreak{}v4\_\allowbreak{}3.md}\par
{\footnotesize \textit{Located.} Sections 9.3-9.5: isotropic cap theorem, exact fixed-side contraction, support audits and tuned comparison}\par
{\footnotesize \textit{Scope.} Separate isotropic SU(3) pentagonal cap model; complete fourth-order connected offsite coefficient modulo scalar, with stated support realizations}\par
{\footnotesize SHA-256 \texttt{1f8295c3a797aeb8} \quad graph: \texttt{CITE:\allowbreak{}UNIFIED}, \texttt{CITE:v4.3}}\par
\smallskip
\noindent\textbf{\texttt{fourthCapRun}} \hfill {\footnotesize verified}\par
{\small Target-blind pentagonal fourth-order run}\par
{\footnotesize\ttfamily C:/\allowbreak{}WORKHOUSE/\allowbreak{}REPO/\allowbreak{}runs/\allowbreak{}blind\_\allowbreak{}pentagonal\_\allowbreak{}o4\_\allowbreak{}2026-\allowbreak{}08-\allowbreak{}28/\allowbreak{}README.md}\par
{\footnotesize \textit{Located.} What it produced; What it changes; pinned exact backend and 24 internal gates}\par
{\footnotesize \textit{Scope.} Retained 2026-08-28 cold run on all 48 fixed-side histories, exact endpoints and Q-chain proper returns}\par
{\footnotesize SHA-256 \texttt{8f89f36a51ead640}}\par
\smallskip
\noindent\textbf{\texttt{fourthHolonomy}} \hfill {\footnotesize verified}\par
{\small Flat-holonomy Wilson sources and a repair for the changing physical rank}\par
{\footnotesize\ttfamily C:/\allowbreak{}WORKHOUSE/\allowbreak{}REPO/\allowbreak{}docs/\allowbreak{}derivations/\allowbreak{}wilson-\allowbreak{}spatial-\allowbreak{}inputs/\allowbreak{}G19\_\allowbreak{}FLAT\_\allowbreak{}HOLONOMY\_\allowbreak{}SOURCES\_\allowbreak{}AND\_\allowbreak{}RANK\_\allowbreak{}REPAIR\_\allowbreak{}20260907.md}\par
{\footnotesize \textit{Located.} Sections 1-3, equations (1)-(13): full flat-background form, twisted cochains, shifted alias grids and physical ranks; Section 4, equations (14)-(16): right inverse, smooth redundant sources and nonflat submersion; Section 5, equations (17)-(18): SU(2) norm-one rank jump; Sections 6-7: interacting and verification scope}\par
{\footnotesize \textit{Scope.} Analytic all-volume theorem on labeled periodic cubic lattices n=Lm>=3, L>=2, all periodic flat compact-unitary-group connections in the original-link norm; actual anchored matrix observation. Nonflat submersion is a source-geometric estimate in the Hilbert-Schmidt metric, not a quantum fast form or source-marginal comparison}\par
{\footnotesize SHA-256 \texttt{c6d524158f83d3a1} \quad graph: \texttt{CITE:\allowbreak{}WILSON\_\allowbreak{}SPATIAL\_\allowbreak{}FLAT\_\allowbreak{}INPUT}}\par
\smallskip
\noindent\textbf{\texttt{fourthHolonomyControls}} \hfill {\footnotesize verified}\par
{\small Flat-background path sources and the physical rank repair: exact controls}\par
{\footnotesize\ttfamily C:/\allowbreak{}WORKHOUSE/\allowbreak{}REPO/\allowbreak{}runs/\allowbreak{}flat\_\allowbreak{}holonomy\_\allowbreak{}sources\_\allowbreak{}2026-\allowbreak{}09-\allowbreak{}07/\allowbreak{}README.md}\par
{\footnotesize \textit{Located.} Reproduction command and exact-control scope: all Fourier/alias directions at n=4,L=2 for neutral and three charged backgrounds; direct real-space boundary controls at (4,2),(3,3),(4,4); projection-jump witness and mathematical negative controls}\par
{\footnotesize \textit{Scope.} Historical exact finite controls recorded in Q(i), with frozen source/checker hashes. No fresh complete replay or new Lean formalization is claimed by the reference-paper edit; the new companion script separately checks stated algebraic constants, small exact witness identities and pinned evidence consistency}\par
{\footnotesize SHA-256 \texttt{2be8d86c821bee99}}\par
\smallskip
\noindent\textbf{\texttt{fourthHolonomyLocal}} \hfill {\footnotesize verified}\par
{\small Local gauge-covariant path averages with a uniform transverse tangent bound}\par
{\footnotesize\ttfamily C:/\allowbreak{}WORKHOUSE/\allowbreak{}REPO/\allowbreak{}docs/\allowbreak{}derivations/\allowbreak{}wilson-\allowbreak{}spatial-\allowbreak{}inputs/\allowbreak{}G19\_\allowbreak{}LOCAL\_\allowbreak{}COVARIANT\_\allowbreak{}AVERAGED\_\allowbreak{}PATH\_\allowbreak{}SOURCES\_\allowbreak{}20260905.md}\par
{\footnotesize \textit{Located.} Section 1 labeled-cell and counting-norm conventions; Section 2 equations (2)-(3) actual anchored path average, endpoint covariance and two-box support; Sections 3-5 equations (4)-(18) distinct zero-cochain maps, alias matching and full form proof}\par
{\footnotesize \textit{Scope.} Globally defined nonlinear matrix average, exact endpoint covariance and local support. Original tangent theorem at identity; the fourthHolonomy source extends it to all flat winding backgrounds. Coordinate-ordered anchors do not by themselves establish rotation or reflection covariance}\par
{\footnotesize SHA-256 \texttt{32f04e9b3682463f} \quad graph: \texttt{CITE:\allowbreak{}SEPT\_\allowbreak{}DOC\_\allowbreak{}G19\_\allowbreak{}LOCAL\_\allowbreak{}COVARIANT\_\allowbreak{}AVERAGED\_\allowbreak{}PATH\_\allowbreak{}SOURCES\_\allowbreak{}20260905}}\par
\smallskip
\noindent\textbf{\texttt{fourthMicroCellular}} \hfill {\footnotesize verified}\par
{\small Universal cellular Hodge-Feshbach algebra and tetrahedral scope}\par
{\footnotesize\ttfamily C:/\allowbreak{}WORKHOUSE/\allowbreak{}REPO/\allowbreak{}docs/\allowbreak{}derivations/\allowbreak{}universal-\allowbreak{}cellular-\allowbreak{}hodge-\allowbreak{}tetrahedral.md}\par
{\footnotesize \textit{Located.} Section 2, Theorems 1-3: connected orientable single-cell boundary identities; Section 3.1-3.2, Theorems 4-5: tetrahedral Hodge matrices and full S4 commutant; Section 3.4: physical-history identification}\par
{\footnotesize \textit{Scope.} Unweighted cellular face chains on a regular cellulation of a sphere; tetrahedral permutation representation. Physical Haar/Fierz identification of an actual projected history remains separate}\par
{\footnotesize SHA-256 \texttt{427974f8eda3ae77} \quad graph: \texttt{CITE:\allowbreak{}UNIVERSAL\_\allowbreak{}CELLULAR\_\allowbreak{}HODGE\_\allowbreak{}TETRAHEDRAL}}\par
\smallskip
\noindent\textbf{\texttt{fourthMicroEngine}} \hfill {\footnotesize verified}\par
{\small Exact SU(3) marked-cluster engine: determinant-family invariant projectors}\par
{\footnotesize\ttfamily C:/\allowbreak{}WORKHOUSE/\allowbreak{}REPO/\allowbreak{}corpus-\allowbreak{}import/\allowbreak{}programs/\allowbreak{}hodge\_\allowbreak{}o4\_\allowbreak{}adjudication/\allowbreak{}src/\allowbreak{}DATA\_\allowbreak{}SU3\_\allowbreak{}Exact\_\allowbreak{}MarkedCluster\_\allowbreak{}m4\_\allowbreak{}Colab.py}\par
{\footnotesize \textit{Located.} Lines 810-823: mixed Gram and pseudoinverse; lines 825-967: ten pure-six tensors, complete Gram contraction, pseudoinverse, coefficient projector, 456 branches and exact gates}\par
{\footnotesize \textit{Scope.} Local SU(3) Haar projectors in the (3,0), (4,1), (6,0) families and conjugates. Pure-six Gram is rank five; a local projector is not a completed excited connected coefficient}\par
{\footnotesize SHA-256 \texttt{be9d77f5b245715e} \quad graph: \texttt{CORPUS:programs-hodge-o4-adjudication-src-data--e640fc}}\par
\smallskip
\noindent\textbf{\texttt{fourthMicroEpsilon}} \hfill {\footnotesize verified}\par
{\small Exact SU(3) Haar integration by epsilon-network reduction}\par
{\footnotesize\ttfamily C:/\allowbreak{}WORKHOUSE/\allowbreak{}REPO/\allowbreak{}src/\allowbreak{}workhouse/\allowbreak{}haar\_\allowbreak{}epsilon.py}\par
{\footnotesize \textit{Located.} Module derivation and projector constructors, lines 1-123; reduce\_\allowbreak{}network, lines 157-235; haar\_\allowbreak{}inner, lines 240-287}\par
{\footnotesize \textit{Scope.} Exact delta-epsilon contraction of the source Haar projectors. Retains all supported determinant families and uses the original balanced-link contraction}\par
{\footnotesize SHA-256 \texttt{4102070900394dc1}}\par
\smallskip
\noindent\textbf{\texttt{fourthMicroEpsilonTests}} \hfill {\footnotesize verified}\par
{\small Independent exact comparisons for the epsilon-network Haar integrator}\par
{\footnotesize\ttfamily C:/\allowbreak{}WORKHOUSE/\allowbreak{}REPO/\allowbreak{}tests/\allowbreak{}test\_\allowbreak{}haar\_\allowbreak{}epsilon.py}\par
{\footnotesize \textit{Located.} Closed forms and source-engine comparisons, lines 38-72; the two fixed pure-six shared-link references, lines 76-106; exact triality orthogonality, lines 109-112}\par
{\footnotesize \textit{Scope.} Bounded local integration comparisons, including two pure-six source references with values +1 and -1. Not a full connected-kernel replay}\par
{\footnotesize SHA-256 \texttt{42d670ffc6307b87}}\par
\smallskip
\noindent\textbf{\texttt{fourthMicroLoopcalc}} \hfill {\footnotesize verified}\par
{\small Exact Wilson-loop calculus with weighted private paths}\par
{\footnotesize\ttfamily C:/\allowbreak{}WORKHOUSE/\allowbreak{}REPO/\allowbreak{}src/\allowbreak{}workhouse/\allowbreak{}loopcalc.py}\par
{\footnotesize \textit{Located.} LINK\_\allowbreak{}WEIGHT, lines 54-64; h0\_\allowbreak{}link and apply\_\allowbreak{}h0, lines 291-337; reduced\_\allowbreak{}words and \_\allowbreak{}flush\_\allowbreak{}run, lines 593-635; projected second-order words, lines 690-699}\par
{\footnotesize \textit{Scope.} Rank-specific exact operator implementation. The paper's general path-isometry proof makes product dependence and compatible retained projections explicit}\par
{\footnotesize SHA-256 \texttt{c707a1701e526741}}\par
\smallskip
\noindent\textbf{\texttt{fourthMicroPathRecord}} \hfill {\footnotesize verified}\par
{\small Exact rational certificate for full and reduced three-face clusters}\par
{\footnotesize\ttfamily C:/\allowbreak{}WORKHOUSE/\allowbreak{}REPO/\allowbreak{}runs/\allowbreak{}path\_\allowbreak{}reduction\_\allowbreak{}2026-\allowbreak{}09-\allowbreak{}04/\allowbreak{}certificate.json}\par
{\footnotesize \textit{Located.} reduced\_\allowbreak{}vs\_\allowbreak{}full: eight cluster types, effective-link weights and equality of every odd/even channel over Q(N); by\_\allowbreak{}history: two-hop and dressing correspondence}\par
{\footnotesize \textit{Scope.} Retained complete eight-cluster symbolic comparison. The new bounded replay independently checks full/reduced physical pair H2 and squared-resolvent words; it does not rerun this entire campaign}\par
{\footnotesize SHA-256 \texttt{8a993b93b2c9f0a6}}\par
\smallskip
\noindent\textbf{\texttt{fourthMicroPaths}} \hfill {\footnotesize verified}\par
{\small Private-link paths are single links, and universality holds history by history}\par
{\footnotesize\ttfamily C:/\allowbreak{}WORKHOUSE/\allowbreak{}REPO/\allowbreak{}docs/\allowbreak{}decisions/\allowbreak{}0030-\allowbreak{}private-\allowbreak{}link-\allowbreak{}paths-\allowbreak{}are-\allowbreak{}single-\allowbreak{}links.md}\par
{\footnotesize \textit{Located.} What was done: path-reduction lemma, history-level correspondence, and corrected time reversal for the single-contact dressing, lines 18-55}\par
{\footnotesize \textit{Scope.} Product-only private path reduction retains normalized Haar measure and k times the single-link Casimir energy. Recorded fixed-order geometry comparisons are separate from arbitrary-multiplicity identities}\par
{\footnotesize SHA-256 \texttt{b9475f43050aed28} \quad graph: \texttt{ADR:0030}}\par
\smallskip
\noindent\textbf{\texttt{g9combined}} \hfill {\footnotesize verified}\par
{\small G9: combined sixth-order support and an unavoidable carrier shape}\par
{\footnotesize\ttfamily C:/\allowbreak{}WORKHOUSE/\allowbreak{}REPO/\allowbreak{}paper/\allowbreak{}research\_\allowbreak{}notes/\allowbreak{}G9\_\allowbreak{}SIXTH\_\allowbreak{}ORDER\_\allowbreak{}COMBINED\_\allowbreak{}20260911.md}\par
{\footnotesize \textit{Located.} Sections 1-7, especially 2 formal folds, 4 mixing, 5 divisor obstruction, 6 one-face}\par
{\footnotesize \textit{Scope.} Full universal F6 support; conditional noncancellation given H3 factorization and local direct H6; exact complete SU3 one-face test, no multiface H6}\par
{\footnotesize SHA-256 \texttt{a48d276398e9e1c3} \quad graph: \texttt{CITE:\allowbreak{}G9\_\allowbreak{}SIXTH\_\allowbreak{}ORDER\_\allowbreak{}COMBINED}}\par
\smallskip
\noindent\textbf{\texttt{g9direct}} \hfill {\footnotesize verified}\par
{\small Direct sixth order on actual clusters over Q(N): one face complete, the pair stages scripted}\par
{\footnotesize\ttfamily C:/\allowbreak{}WORKHOUSE/\allowbreak{}REPO/\allowbreak{}runs/\allowbreak{}g9\_\allowbreak{}direct\_\allowbreak{}h6\_\allowbreak{}pair\_\allowbreak{}2026-\allowbreak{}09-\allowbreak{}11/\allowbreak{}README.md}\par
{\footnotesize \textit{Located.} What is computed stage table; Result; Where it stops}\par
{\footnotesize \textit{Scope.} One-face Q(N) H6 complete and min\_\allowbreak{}rank9; both pair H6 stages unexecuted; pair H4 validated}\par
{\footnotesize SHA-256 \texttt{5b0869782f770ca3}}\par
\smallskip
\noindent\textbf{\texttt{guideAgreement}} \hfill {\footnotesize verified}\par
{\small WORKHOUSE working agreement}\par
{\footnotesize\ttfamily C:/\allowbreak{}WORKHOUSE/\allowbreak{}REPO/\allowbreak{}CLAUDE.md}\par
{\footnotesize \textit{Located.} Evidence axes; exact calculations; preservation and scientific source rules}\par
{\footnotesize \textit{Scope.} Operational source; no scientific theorem is inferred from the instructions}\par
{\footnotesize SHA-256 \texttt{639599caf4208186}}\par
\smallskip
\noindent\textbf{\texttt{guideCoordination}} \hfill {\footnotesize verified}\par
{\small Workspace coordination}\par
{\footnotesize\ttfamily C:/\allowbreak{}WORKHOUSE/\allowbreak{}REPO/\allowbreak{}docs/\allowbreak{}workspace\_\allowbreak{}coordination.md}\par
{\footnotesize \textit{Located.} Canonical checkout, preserved pending checkout and runtime routing}\par
{\footnotesize \textit{Scope.} Operational source; no scientific theorem is inferred from the instructions}\par
{\footnotesize SHA-256 \texttt{62160940cdf07354}}\par
\smallskip
\noindent\textbf{\texttt{guideCurrent}} \hfill {\footnotesize changed since curated}\par
{\small Current research map}\par
{\footnotesize\ttfamily C:/\allowbreak{}WORKHOUSE/\allowbreak{}REPO/\allowbreak{}docs/\allowbreak{}current\_\allowbreak{}research.md}\par
{\footnotesize \textit{Located.} Opening current map; Read the current claim at its own scope; Select the kind of work; Keep the corrections attached; Local class spectrum continuation}\par
{\footnotesize \textit{Scope.} Maintained source-level goal/current navigation. Operative mathematical claims must be read with their exact linked derivation; historical next-step prose is not promoted to a live blocker}\par
{\footnotesize SHA-256 \texttt{c326b036011e43bf} \quad \textbf{now 60331606c49b32c3}}\par
\smallskip
\noindent\textbf{\texttt{guideDiscovery}} \hfill {\footnotesize verified}\par
{\small Corpus discovery for research agents}\par
{\footnotesize\ttfamily C:/\allowbreak{}WORKHOUSE/\allowbreak{}REPO/\allowbreak{}docs/\allowbreak{}graph\_\allowbreak{}discovery.md}\par
{\footnotesize \textit{Located.} Search scope, query fusion, path-pair dossiers and source-reading boundary}\par
{\footnotesize \textit{Scope.} Operational source; no scientific theorem is inferred from the instructions}\par
{\footnotesize SHA-256 \texttt{9c0487b01473d252}}\par
\smallskip
\noindent\textbf{\texttt{guideFormalization}} \hfill {\footnotesize verified}\par
{\small Formalization and proof connections}\par
{\footnotesize\ttfamily C:/\allowbreak{}WORKHOUSE/\allowbreak{}REPO/\allowbreak{}docs/\allowbreak{}formalization\_\allowbreak{}workflow.md}\par
{\footnotesize \textit{Located.} Source records, whole versus partial formal coverage, exact hashes and integration order}\par
{\footnotesize \textit{Scope.} Operational source; no scientific theorem is inferred from the instructions}\par
{\footnotesize SHA-256 \texttt{f6bdf26655b410e5}}\par
\smallskip
\noindent\textbf{\texttt{guideGoal}} \hfill {\footnotesize verified}\par
{\small Research goal and the path to it}\par
{\footnotesize\ttfamily C:/\allowbreak{}WORKHOUSE/\allowbreak{}REPO/\allowbreak{}docs/\allowbreak{}research\_\allowbreak{}goal.md}\par
{\footnotesize \textit{Located.} Governing objective; Established inputs keep their scopes; Remaining obligations for the continuum target; Formalization is a separate dependency chain}\par
{\footnotesize \textit{Scope.} Maintained source-level goal/current navigation. Operative mathematical claims must be read with their exact linked derivation; historical next-step prose is not promoted to a live blocker}\par
{\footnotesize SHA-256 \texttt{3827fabfb22a821c}}\par
\smallskip
\noindent\textbf{\texttt{guideGraph091d393722}} \hfill {\footnotesize verified}\par
{\small Theory-graph crosswalk input: lean\_\allowbreak{}dependencies.json}\par
{\footnotesize\ttfamily C:/\allowbreak{}WORKHOUSE/\allowbreak{}REPO/\allowbreak{}ledger/\allowbreak{}lean\_\allowbreak{}dependencies.json}\par
{\footnotesize \textit{Located.} Exact ID, evidence-symbol and source-locator validation for the 36-statement crosswalk}\par
{\footnotesize \textit{Scope.} Registry or scoped evidence input, not an additional scientific theorem or a fresh run}\par
{\footnotesize SHA-256 \texttt{b44bcc06e7c3429b}}\par
\smallskip
\noindent\textbf{\texttt{guideGraph09e1a91be2}} \hfill {\footnotesize verified}\par
{\small Theory-graph crosswalk input: theorems.yaml}\par
{\footnotesize\ttfamily C:/\allowbreak{}WORKHOUSE/\allowbreak{}REPO/\allowbreak{}ledger/\allowbreak{}theorems.yaml}\par
{\footnotesize \textit{Located.} Exact ID, evidence-symbol and source-locator validation for the 36-statement crosswalk}\par
{\footnotesize \textit{Scope.} Registry or scoped evidence input, not an additional scientific theorem or a fresh run}\par
{\footnotesize SHA-256 \texttt{52fa82a10e1df02e}}\par
\smallskip
\noindent\textbf{\texttt{guideGraph0fb13070b1}} \hfill {\footnotesize verified}\par
{\small Theory-graph crosswalk input: odd\_\allowbreak{}order.py}\par
{\footnotesize\ttfamily C:/\allowbreak{}WORKHOUSE/\allowbreak{}REPO/\allowbreak{}src/\allowbreak{}workhouse/\allowbreak{}invariants/\allowbreak{}odd\_\allowbreak{}order.py}\par
{\footnotesize \textit{Located.} Exact ID, evidence-symbol and source-locator validation for the 36-statement crosswalk}\par
{\footnotesize \textit{Scope.} Registry or scoped evidence input, not an additional scientific theorem or a fresh run}\par
{\footnotesize SHA-256 \texttt{1e0371f0ece14cd5}}\par
\smallskip
\noindent\textbf{\texttt{guideGraph1f7afb2791}} \hfill {\footnotesize verified}\par
{\small Theory-graph crosswalk input: wilson\_\allowbreak{}marked.py}\par
{\footnotesize\ttfamily C:/\allowbreak{}WORKHOUSE/\allowbreak{}REPO/\allowbreak{}src/\allowbreak{}workhouse/\allowbreak{}invariants/\allowbreak{}wilson\_\allowbreak{}marked.py}\par
{\footnotesize \textit{Located.} Exact ID, evidence-symbol and source-locator validation for the 36-statement crosswalk}\par
{\footnotesize \textit{Scope.} Registry or scoped evidence input, not an additional scientific theorem or a fresh run}\par
{\footnotesize SHA-256 \texttt{e28c04613e93f4a6}}\par
\smallskip
\noindent\textbf{\texttt{guideGraph2b74998cc0}} \hfill {\footnotesize changed since curated}\par
{\small Theory-graph crosswalk input: claims.jsonl}\par
{\footnotesize\ttfamily C:/\allowbreak{}WORKHOUSE/\allowbreak{}REPO/\allowbreak{}index/\allowbreak{}claims.jsonl}\par
{\footnotesize \textit{Located.} Exact ID, evidence-symbol and source-locator validation for the 36-statement crosswalk}\par
{\footnotesize \textit{Scope.} Registry or scoped evidence input, not an additional scientific theorem or a fresh run}\par
{\footnotesize SHA-256 \texttt{de7f4d2cf88a7c93} \quad \textbf{now 13fe8385396e46fd}}\par
\smallskip
\noindent\textbf{\texttt{guideGraph320d0a5fa5}} \hfill {\footnotesize verified}\par
{\small Theory-graph crosswalk input: local\_\allowbreak{}class\_\allowbreak{}wick.py}\par
{\footnotesize\ttfamily C:/\allowbreak{}WORKHOUSE/\allowbreak{}REPO/\allowbreak{}src/\allowbreak{}workhouse/\allowbreak{}invariants/\allowbreak{}local\_\allowbreak{}class\_\allowbreak{}wick.py}\par
{\footnotesize \textit{Located.} Exact ID, evidence-symbol and source-locator validation for the 36-statement crosswalk}\par
{\footnotesize \textit{Scope.} Registry or scoped evidence input, not an additional scientific theorem or a fresh run}\par
{\footnotesize SHA-256 \texttt{e5a330ca8a202bfd}}\par
\smallskip
\noindent\textbf{\texttt{guideGraph3ad8b3a7d5}} \hfill {\footnotesize verified}\par
{\small Theory-graph crosswalk input: moving\_\allowbreak{}time\_\allowbreak{}gap.py}\par
{\footnotesize\ttfamily C:/\allowbreak{}WORKHOUSE/\allowbreak{}REPO/\allowbreak{}src/\allowbreak{}workhouse/\allowbreak{}invariants/\allowbreak{}moving\_\allowbreak{}time\_\allowbreak{}gap.py}\par
{\footnotesize \textit{Located.} Exact ID, evidence-symbol and source-locator validation for the 36-statement crosswalk}\par
{\footnotesize \textit{Scope.} Registry or scoped evidence input, not an additional scientific theorem or a fresh run}\par
{\footnotesize SHA-256 \texttt{9aa42e2893e59834}}\par
\smallskip
\noindent\textbf{\texttt{guideGraph3d6bc5d9b5}} \hfill {\footnotesize verified}\par
{\small Theory-graph crosswalk input: AnisotropyVariance.lean}\par
{\footnotesize\ttfamily C:/\allowbreak{}WORKHOUSE/\allowbreak{}REPO/\allowbreak{}lean/\allowbreak{}Workhouse/\allowbreak{}AnisotropyVariance.lean}\par
{\footnotesize \textit{Located.} Exact ID, evidence-symbol and source-locator validation for the 36-statement crosswalk}\par
{\footnotesize \textit{Scope.} Registry or scoped evidence input, not an additional scientific theorem or a fresh run}\par
{\footnotesize SHA-256 \texttt{8d7c8fc2006c3457}}\par
\smallskip
\noindent\textbf{\texttt{guideGraph44ed90e13c}} \hfill {\footnotesize verified}\par
{\small Theory-graph crosswalk input: sixth\_\allowbreak{}order.py}\par
{\footnotesize\ttfamily C:/\allowbreak{}WORKHOUSE/\allowbreak{}REPO/\allowbreak{}src/\allowbreak{}workhouse/\allowbreak{}invariants/\allowbreak{}sixth\_\allowbreak{}order.py}\par
{\footnotesize \textit{Located.} Exact ID, evidence-symbol and source-locator validation for the 36-statement crosswalk}\par
{\footnotesize \textit{Scope.} Registry or scoped evidence input, not an additional scientific theorem or a fresh run}\par
{\footnotesize SHA-256 \texttt{7ddde397d7297f52}}\par
\smallskip
\noindent\textbf{\texttt{guideGraph4bfd0509d6}} \hfill {\footnotesize verified}\par
{\small Theory-graph crosswalk input: derivation\_\allowbreak{}statements.yaml}\par
{\footnotesize\ttfamily C:/\allowbreak{}WORKHOUSE/\allowbreak{}REPO/\allowbreak{}ledger/\allowbreak{}derivation\_\allowbreak{}statements.yaml}\par
{\footnotesize \textit{Located.} Exact ID, evidence-symbol and source-locator validation for the 36-statement crosswalk}\par
{\footnotesize \textit{Scope.} Registry or scoped evidence input, not an additional scientific theorem or a fresh run}\par
{\footnotesize SHA-256 \texttt{2591dd0e9447e68d}}\par
\smallskip
\noindent\textbf{\texttt{guideGraph4da7303fcd}} \hfill {\footnotesize verified}\par
{\small Theory-graph crosswalk input: archive\_\allowbreak{}derivations.py}\par
{\footnotesize\ttfamily C:/\allowbreak{}WORKHOUSE/\allowbreak{}REPO/\allowbreak{}src/\allowbreak{}workhouse/\allowbreak{}invariants/\allowbreak{}archive\_\allowbreak{}derivations.py}\par
{\footnotesize \textit{Located.} Exact ID, evidence-symbol and source-locator validation for the 36-statement crosswalk}\par
{\footnotesize \textit{Scope.} Registry or scoped evidence input, not an additional scientific theorem or a fresh run}\par
{\footnotesize SHA-256 \texttt{86ddd481032a5a46}}\par
\smallskip
\noindent\textbf{\texttt{guideGraph4e7cccef42}} \hfill {\footnotesize verified}\par
{\small Theory-graph crosswalk input: Basic.lean}\par
{\footnotesize\ttfamily C:/\allowbreak{}WORKHOUSE/\allowbreak{}REPO/\allowbreak{}lean/\allowbreak{}Workhouse/\allowbreak{}Basic.lean}\par
{\footnotesize \textit{Located.} Exact ID, evidence-symbol and source-locator validation for the 36-statement crosswalk}\par
{\footnotesize \textit{Scope.} Registry or scoped evidence input, not an additional scientific theorem or a fresh run}\par
{\footnotesize SHA-256 \texttt{0ba5aa78230ed19a}}\par
\smallskip
\noindent\textbf{\texttt{guideGraph547cb9b68b}} \hfill {\footnotesize verified}\par
{\small Theory-graph crosswalk input: briefing.py}\par
{\footnotesize\ttfamily C:/\allowbreak{}WORKHOUSE/\allowbreak{}REPO/\allowbreak{}src/\allowbreak{}workhouse/\allowbreak{}briefing.py}\par
{\footnotesize \textit{Located.} Exact ID, evidence-symbol and source-locator validation for the 36-statement crosswalk}\par
{\footnotesize \textit{Scope.} Registry or scoped evidence input, not an additional scientific theorem or a fresh run}\par
{\footnotesize SHA-256 \texttt{6eee69c2369cc382}}\par
\smallskip
\noindent\textbf{\texttt{guideGraph5981cf40fb}} \hfill {\footnotesize verified}\par
{\small Theory-graph crosswalk input: ThermodynamicLimit.lean}\par
{\footnotesize\ttfamily C:/\allowbreak{}WORKHOUSE/\allowbreak{}REPO/\allowbreak{}lean/\allowbreak{}Workhouse/\allowbreak{}ThermodynamicLimit.lean}\par
{\footnotesize \textit{Located.} Exact ID, evidence-symbol and source-locator validation for the 36-statement crosswalk}\par
{\footnotesize \textit{Scope.} Registry or scoped evidence input, not an additional scientific theorem or a fresh run}\par
{\footnotesize SHA-256 \texttt{b5c50685452c1f58}}\par
\smallskip
\noindent\textbf{\texttt{guideGraph6d087a1389}} \hfill {\footnotesize verified}\par
{\small Theory-graph crosswalk input: TheoryCurrentBridges.lean}\par
{\footnotesize\ttfamily C:/\allowbreak{}WORKHOUSE/\allowbreak{}REPO/\allowbreak{}lean/\allowbreak{}Workhouse/\allowbreak{}TheoryCurrentBridges.lean}\par
{\footnotesize \textit{Located.} Exact ID, evidence-symbol and source-locator validation for the 36-statement crosswalk}\par
{\footnotesize \textit{Scope.} Registry or scoped evidence input, not an additional scientific theorem or a fresh run}\par
{\footnotesize SHA-256 \texttt{52c4309513867c0f}}\par
\smallskip
\noindent\textbf{\texttt{guideGraph756bf40850}} \hfill {\footnotesize verified}\par
{\small Theory-graph crosswalk input: GroundStateAssembly.lean}\par
{\footnotesize\ttfamily C:/\allowbreak{}WORKHOUSE/\allowbreak{}REPO/\allowbreak{}lean/\allowbreak{}Workhouse/\allowbreak{}GroundStateAssembly.lean}\par
{\footnotesize \textit{Located.} Exact ID, evidence-symbol and source-locator validation for the 36-statement crosswalk}\par
{\footnotesize \textit{Scope.} Registry or scoped evidence input, not an additional scientific theorem or a fresh run}\par
{\footnotesize SHA-256 \texttt{4e5474836257d0ee}}\par
\smallskip
\noindent\textbf{\texttt{guideGraph75f6911b04}} \hfill {\footnotesize verified}\par
{\small Theory-graph crosswalk input: theory\_\allowbreak{}current\_\allowbreak{}bridges.py}\par
{\footnotesize\ttfamily C:/\allowbreak{}WORKHOUSE/\allowbreak{}REPO/\allowbreak{}src/\allowbreak{}workhouse/\allowbreak{}invariants/\allowbreak{}theory\_\allowbreak{}current\_\allowbreak{}bridges.py}\par
{\footnotesize \textit{Located.} Exact ID, evidence-symbol and source-locator validation for the 36-statement crosswalk}\par
{\footnotesize \textit{Scope.} Registry or scoped evidence input, not an additional scientific theorem or a fresh run}\par
{\footnotesize SHA-256 \texttt{1076d69d4df25e3d}}\par
\smallskip
\noindent\textbf{\texttt{guideGraph848c5f0441}} \hfill {\footnotesize verified}\par
{\small Theory-graph crosswalk input: ResolventLocalization.lean}\par
{\footnotesize\ttfamily C:/\allowbreak{}WORKHOUSE/\allowbreak{}REPO/\allowbreak{}lean/\allowbreak{}Workhouse/\allowbreak{}ResolventLocalization.lean}\par
{\footnotesize \textit{Located.} Exact ID, evidence-symbol and source-locator validation for the 36-statement crosswalk}\par
{\footnotesize \textit{Scope.} Registry or scoped evidence input, not an additional scientific theorem or a fresh run}\par
{\footnotesize SHA-256 \texttt{f0a3d97ec2d760cc}}\par
\smallskip
\noindent\textbf{\texttt{guideGraph899902d358}} \hfill {\footnotesize verified}\par
{\small Theory-graph crosswalk input: rank\_\allowbreak{}law.py}\par
{\footnotesize\ttfamily C:/\allowbreak{}WORKHOUSE/\allowbreak{}REPO/\allowbreak{}src/\allowbreak{}workhouse/\allowbreak{}invariants/\allowbreak{}rank\_\allowbreak{}law.py}\par
{\footnotesize \textit{Located.} Exact ID, evidence-symbol and source-locator validation for the 36-statement crosswalk}\par
{\footnotesize \textit{Scope.} Registry or scoped evidence input, not an additional scientific theorem or a fresh run}\par
{\footnotesize SHA-256 \texttt{ff3d2c9f580d1227}}\par
\smallskip
\noindent\textbf{\texttt{guideGraph8eb2974983}} \hfill {\footnotesize verified}\par
{\small Theory-graph crosswalk input: theory\_\allowbreak{}graph\_\allowbreak{}protocol.md}\par
{\footnotesize\ttfamily C:/\allowbreak{}WORKHOUSE/\allowbreak{}REPO/\allowbreak{}docs/\allowbreak{}theory\_\allowbreak{}graph\_\allowbreak{}protocol.md}\par
{\footnotesize \textit{Located.} Exact ID, evidence-symbol and source-locator validation for the 36-statement crosswalk}\par
{\footnotesize \textit{Scope.} Registry or scoped evidence input, not an additional scientific theorem or a fresh run}\par
{\footnotesize SHA-256 \texttt{54cbbe206808fa30}}\par
\smallskip
\noindent\textbf{\texttt{guideGraph9de14e9119}} \hfill {\footnotesize verified}\par
{\small Theory-graph crosswalk input: os\_\allowbreak{}kernel\_\allowbreak{}gap.py}\par
{\footnotesize\ttfamily C:/\allowbreak{}WORKHOUSE/\allowbreak{}REPO/\allowbreak{}src/\allowbreak{}workhouse/\allowbreak{}invariants/\allowbreak{}os\_\allowbreak{}kernel\_\allowbreak{}gap.py}\par
{\footnotesize \textit{Located.} Exact ID, evidence-symbol and source-locator validation for the 36-statement crosswalk}\par
{\footnotesize \textit{Scope.} Registry or scoped evidence input, not an additional scientific theorem or a fresh run}\par
{\footnotesize SHA-256 \texttt{19e4bfcd0fdd60db}}\par
\smallskip
\noindent\textbf{\texttt{guideGrapha62e6f677b}} \hfill {\footnotesize verified}\par
{\small Theory-graph crosswalk input: spectral\_\allowbreak{}budgets.py}\par
{\footnotesize\ttfamily C:/\allowbreak{}WORKHOUSE/\allowbreak{}REPO/\allowbreak{}src/\allowbreak{}workhouse/\allowbreak{}invariants/\allowbreak{}spectral\_\allowbreak{}budgets.py}\par
{\footnotesize \textit{Located.} Exact ID, evidence-symbol and source-locator validation for the 36-statement crosswalk}\par
{\footnotesize \textit{Scope.} Registry or scoped evidence input, not an additional scientific theorem or a fresh run}\par
{\footnotesize SHA-256 \texttt{5632202088a3397d}}\par
\smallskip
\noindent\textbf{\texttt{guideGrapha7010edcef}} \hfill {\footnotesize verified}\par
{\small Theory-graph crosswalk input: wilson\_\allowbreak{}physical\_\allowbreak{}band.py}\par
{\footnotesize\ttfamily C:/\allowbreak{}WORKHOUSE/\allowbreak{}REPO/\allowbreak{}src/\allowbreak{}workhouse/\allowbreak{}invariants/\allowbreak{}wilson\_\allowbreak{}physical\_\allowbreak{}band.py}\par
{\footnotesize \textit{Located.} Exact ID, evidence-symbol and source-locator validation for the 36-statement crosswalk}\par
{\footnotesize \textit{Scope.} Registry or scoped evidence input, not an additional scientific theorem or a fresh run}\par
{\footnotesize SHA-256 \texttt{e07e351b12946b4f}}\par
\smallskip
\noindent\textbf{\texttt{guideGraphab6984628e}} \hfill {\footnotesize verified}\par
{\small Theory-graph crosswalk input: derivation\_\allowbreak{}formalization.md}\par
{\footnotesize\ttfamily C:/\allowbreak{}WORKHOUSE/\allowbreak{}REPO/\allowbreak{}docs/\allowbreak{}derivation\_\allowbreak{}formalization.md}\par
{\footnotesize \textit{Located.} Exact ID, evidence-symbol and source-locator validation for the 36-statement crosswalk}\par
{\footnotesize \textit{Scope.} Registry or scoped evidence input, not an additional scientific theorem or a fresh run}\par
{\footnotesize SHA-256 \texttt{ab7a04278a58175d}}\par
\smallskip
\noindent\textbf{\texttt{guideGraphaf74744fc5}} \hfill {\footnotesize verified}\par
{\small Theory-graph crosswalk input: HodgeFeshbach.lean}\par
{\footnotesize\ttfamily C:/\allowbreak{}WORKHOUSE/\allowbreak{}REPO/\allowbreak{}lean/\allowbreak{}Workhouse/\allowbreak{}HodgeFeshbach.lean}\par
{\footnotesize \textit{Located.} Exact ID, evidence-symbol and source-locator validation for the 36-statement crosswalk}\par
{\footnotesize \textit{Scope.} Registry or scoped evidence input, not an additional scientific theorem or a fresh run}\par
{\footnotesize SHA-256 \texttt{da3357565a79b1e2}}\par
\smallskip
\noindent\textbf{\texttt{guideGraphc158d5e678}} \hfill {\footnotesize verified}\par
{\small Theory-graph crosswalk input: universal\_\allowbreak{}cellular\_\allowbreak{}hodge.py}\par
{\footnotesize\ttfamily C:/\allowbreak{}WORKHOUSE/\allowbreak{}REPO/\allowbreak{}src/\allowbreak{}workhouse/\allowbreak{}invariants/\allowbreak{}universal\_\allowbreak{}cellular\_\allowbreak{}hodge.py}\par
{\footnotesize \textit{Located.} Exact ID, evidence-symbol and source-locator validation for the 36-statement crosswalk}\par
{\footnotesize \textit{Scope.} Registry or scoped evidence input, not an additional scientific theorem or a fresh run}\par
{\footnotesize SHA-256 \texttt{e5b1d241aafb1a74}}\par
\smallskip
\noindent\textbf{\texttt{guideGraphc54eee0bb2}} \hfill {\footnotesize verified}\par
{\small Theory-graph crosswalk input: README.md}\par
{\footnotesize\ttfamily C:/\allowbreak{}WORKHOUSE/\allowbreak{}REPO/\allowbreak{}runs/\allowbreak{}wilson\_\allowbreak{}physical\_\allowbreak{}band\_\allowbreak{}2026-\allowbreak{}09-\allowbreak{}05/\allowbreak{}README.md}\par
{\footnotesize \textit{Located.} Exact ID, evidence-symbol and source-locator validation for the 36-statement crosswalk}\par
{\footnotesize \textit{Scope.} Registry or scoped evidence input, not an additional scientific theorem or a fresh run}\par
{\footnotesize SHA-256 \texttt{19d829b1e0294415}}\par
\smallskip
\noindent\textbf{\texttt{guideGraphcf78243227}} \hfill {\footnotesize verified}\par
{\small Theory-graph crosswalk input: wilson\_\allowbreak{}shell.py}\par
{\footnotesize\ttfamily C:/\allowbreak{}WORKHOUSE/\allowbreak{}REPO/\allowbreak{}src/\allowbreak{}workhouse/\allowbreak{}invariants/\allowbreak{}wilson\_\allowbreak{}shell.py}\par
{\footnotesize \textit{Located.} Exact ID, evidence-symbol and source-locator validation for the 36-statement crosswalk}\par
{\footnotesize \textit{Scope.} Registry or scoped evidence input, not an additional scientific theorem or a fresh run}\par
{\footnotesize SHA-256 \texttt{d195314d4acc95be}}\par
\smallskip
\noindent\textbf{\texttt{guideGraphd5f378773f}} \hfill {\footnotesize changed since curated}\par
{\small Theory-graph crosswalk input: results.yaml}\par
{\footnotesize\ttfamily C:/\allowbreak{}WORKHOUSE/\allowbreak{}REPO/\allowbreak{}ledger/\allowbreak{}results.yaml}\par
{\footnotesize \textit{Located.} Exact ID, evidence-symbol and source-locator validation for the 36-statement crosswalk}\par
{\footnotesize \textit{Scope.} Registry or scoped evidence input, not an additional scientific theorem or a fresh run}\par
{\footnotesize SHA-256 \texttt{0e0d7bb55ecb2e2c} \quad \textbf{now e9459205198060e1}}\par
\smallskip
\noindent\textbf{\texttt{guideGraphd6828a9cb0}} \hfill {\footnotesize changed since curated}\par
{\small Theory-graph crosswalk input: graph.jsonl}\par
{\footnotesize\ttfamily C:/\allowbreak{}WORKHOUSE/\allowbreak{}REPO/\allowbreak{}index/\allowbreak{}graph.jsonl}\par
{\footnotesize \textit{Located.} Exact ID, evidence-symbol and source-locator validation for the 36-statement crosswalk}\par
{\footnotesize \textit{Scope.} Registry or scoped evidence input, not an additional scientific theorem or a fresh run}\par
{\footnotesize SHA-256 \texttt{2e01c8b80bccae21} \quad \textbf{now 5e42701057bb2e9c}}\par
\smallskip
\noindent\textbf{\texttt{guideGraphe31cf24064}} \hfill {\footnotesize verified}\par
{\small Theory-graph crosswalk input: wilson\_\allowbreak{}spatial.py}\par
{\footnotesize\ttfamily C:/\allowbreak{}WORKHOUSE/\allowbreak{}REPO/\allowbreak{}src/\allowbreak{}workhouse/\allowbreak{}invariants/\allowbreak{}wilson\_\allowbreak{}spatial.py}\par
{\footnotesize \textit{Located.} Exact ID, evidence-symbol and source-locator validation for the 36-statement crosswalk}\par
{\footnotesize \textit{Scope.} Registry or scoped evidence input, not an additional scientific theorem or a fresh run}\par
{\footnotesize SHA-256 \texttt{db3fe6402f39bf07}}\par
\smallskip
\noindent\textbf{\texttt{guideGraphe7cb95cba7}} \hfill {\footnotesize verified}\par
{\small Theory-graph crosswalk input: hodge\_\allowbreak{}feshbach.py}\par
{\footnotesize\ttfamily C:/\allowbreak{}WORKHOUSE/\allowbreak{}REPO/\allowbreak{}src/\allowbreak{}workhouse/\allowbreak{}invariants/\allowbreak{}hodge\_\allowbreak{}feshbach.py}\par
{\footnotesize \textit{Located.} Exact ID, evidence-symbol and source-locator validation for the 36-statement crosswalk}\par
{\footnotesize \textit{Scope.} Registry or scoped evidence input, not an additional scientific theorem or a fresh run}\par
{\footnotesize SHA-256 \texttt{10c8c5014ab4e774}}\par
\smallskip
\noindent\textbf{\texttt{guideGraphff329da461}} \hfill {\footnotesize verified}\par
{\small Theory-graph crosswalk input: path\_\allowbreak{}reduction.py}\par
{\footnotesize\ttfamily C:/\allowbreak{}WORKHOUSE/\allowbreak{}REPO/\allowbreak{}src/\allowbreak{}workhouse/\allowbreak{}invariants/\allowbreak{}path\_\allowbreak{}reduction.py}\par
{\footnotesize \textit{Located.} Exact ID, evidence-symbol and source-locator validation for the 36-statement crosswalk}\par
{\footnotesize \textit{Scope.} Registry or scoped evidence input, not an additional scientific theorem or a fresh run}\par
{\footnotesize SHA-256 \texttt{bb6c58f337ce2ea8}}\par
\smallskip
\noindent\textbf{\texttt{guideLitCatalog}} \hfill {\footnotesize verified}\par
{\small Historical acquisition catalogue, first collection}\par
{\footnotesize\ttfamily C:/\allowbreak{}WORKHOUSE/\allowbreak{}REPO/\allowbreak{}literature/\allowbreak{}library/\allowbreak{}paper\_\allowbreak{}manifest.csv}\par
{\footnotesize \textit{Located.} 99 rows; exact arXiv IDs 2204.12737, 2105.02058, 0704.3244, 1105.0675}\par
{\footnotesize \textit{Scope.} Identity and abstract metadata only; source PDFs not opened or reverified}\par
{\footnotesize SHA-256 \texttt{3c399494a6e04724}}\par
\smallskip
\noindent\textbf{\texttt{guideLitCatalogTwo}} \hfill {\footnotesize verified}\par
{\small Historical acquisition catalogue, second collection}\par
{\footnotesize\ttfamily C:/\allowbreak{}WORKHOUSE/\allowbreak{}REPO/\allowbreak{}literature/\allowbreak{}library/\allowbreak{}volume2/\allowbreak{}paper\_\allowbreak{}manifest.csv}\par
{\footnotesize \textit{Located.} 161 rows; discovery\_\allowbreak{}path and workhouse\_\allowbreak{}connection fields}\par
{\footnotesize \textit{Scope.} Acquisition metadata only}\par
{\footnotesize SHA-256 \texttt{2f6b178eb74924e8}}\par
\smallskip
\noindent\textbf{\texttt{guideLitCurvature}} \hfill {\footnotesize verified}\par
{\small Weighted curvature, trial ground states, and a residual gap certificate}\par
{\footnotesize\ttfamily C:/\allowbreak{}WORKHOUSE/\allowbreak{}REPO/\allowbreak{}docs/\allowbreak{}derivations/\allowbreak{}yangmills-\allowbreak{}weighted-\allowbreak{}curvature.md}\par
{\footnotesize \textit{Located.} GST1-GST3 and GST6; source comparison and domain hypotheses}\par
{\footnotesize \textit{Scope.} Local differential, form and spectral derivation; conditional global curvature/residual application}\par
{\footnotesize SHA-256 \texttt{75c853725c68db30} \quad graph: \texttt{CITE:\allowbreak{}YM\_\allowbreak{}GROUND\_\allowbreak{}STATE}}\par
\smallskip
\noindent\textbf{\texttt{guideLitIndex}} \hfill {\footnotesize verified}\par
{\small Native published-literature evidence map}\par
{\footnotesize\ttfamily C:/\allowbreak{}WORKHOUSE/\allowbreak{}REPO/\allowbreak{}literature/\allowbreak{}index.yaml}\par
{\footnotesize \textit{Located.} papers: selected IDs; bears\_\allowbreak{}on edges and reading-status notes}\par
{\footnotesize \textit{Scope.} Native evidence metadata, distinct from this audit reading status}\par
{\footnotesize SHA-256 \texttt{6880e1ec83b98068} \quad graph: \texttt{LIT:\allowbreak{}AIZENMAN\_\allowbreak{}1982}, \texttt{LIT:\allowbreak{}AIZENMAN\_\allowbreak{}1982:\allowbreak{}G19}, \texttt{LIT:\allowbreak{}AT\_\allowbreak{}2020\_\allowbreak{}SU3}, \texttt{LIT:\allowbreak{}AT\_\allowbreak{}2020\_\allowbreak{}SU3:\allowbreak{}G18}, \texttt{LIT:\allowbreak{}AT\_\allowbreak{}2021\_\allowbreak{}SUN}, \texttt{LIT:\allowbreak{}AT\_\allowbreak{}2021\_\allowbreak{}SUN:\allowbreak{}G19}}\par
\smallskip
\noindent\textbf{\texttt{guideLitJW}} \hfill {\footnotesize outside this checkout}\par
{\small Quantum Yang-Mills Theory}\par
{\footnotesize\ttfamily C:/\allowbreak{}WORKHOUSE/\allowbreak{}ALL THEORY/\allowbreak{}WORKHOUSE/\allowbreak{}literature/\allowbreak{}inbox/\allowbreak{}JW\_\allowbreak{}2006/\allowbreak{}extracted/\allowbreak{}JW\_\allowbreak{}2006.txt}\par
{\footnotesize \textit{Located.} PDF pages 5-6, sections 3-5; pages 11-12, sections 6.5-6.6}\par
{\footnotesize \textit{Scope.} External source text; selected passages actually read; local applications retain their own proof hypotheses}\par
{\footnotesize SHA-256 \texttt{dab56d50771697fe} \quad reading copy: C:/WORKHOUSE/REPO/literature/inbox/JW\_\allowbreak{}2006/extracted/JW\_\allowbreak{}2006.txt. The extracted-literature guide records its resolution \quad \textbf{workspace path under ALL THEORY/}}\par
\smallskip
\noindent\textbf{\texttt{guideLitMMM}} \hfill {\footnotesize outside this checkout}\par
{\small Orbit Space Curvature as a Source of Mass in Quantum Gauge Theory}\par
{\footnotesize\ttfamily C:/\allowbreak{}WORKHOUSE/\allowbreak{}ALL THEORY/\allowbreak{}WORKHOUSE/\allowbreak{}literature/\allowbreak{}inbox/\allowbreak{}JW\_\allowbreak{}2006/\allowbreak{}extracted/\allowbreak{}MMM\_\allowbreak{}2019\_\allowbreak{}CURVATURE.txt}\par
{\footnotesize \textit{Located.} PDF pages 8-9, II.2-II.9; pages 11-12, II.15-II.23; page 14, II.32-II.36; selected passages only}\par
{\footnotesize \textit{Scope.} External source text; selected passages actually read; local applications retain their own proof hypotheses}\par
{\footnotesize SHA-256 \texttt{3315bb91c1ecac5c} \quad reading copy: C:/WORKHOUSE/REPO/literature/inbox/JW\_\allowbreak{}2006/extracted/MMM\_\allowbreak{}2019\_\allowbreak{}CURVATURE.txt. The extracted-literature guide records its resolution \quad \textbf{workspace path under ALL THEORY/}}\par
\smallskip
\noindent\textbf{\texttt{guideLitMondal}} \hfill {\footnotesize outside this checkout}\par
{\small A Geometric Approach to the Yang-Mills Mass Gap}\par
{\footnotesize\ttfamily C:/\allowbreak{}WORKHOUSE/\allowbreak{}ALL THEORY/\allowbreak{}WORKHOUSE/\allowbreak{}literature/\allowbreak{}inbox/\allowbreak{}JW\_\allowbreak{}2006/\allowbreak{}extracted/\allowbreak{}MONDAL\_\allowbreak{}2023\_\allowbreak{}GEOMETRIC.txt}\par
{\footnotesize \textit{Located.} PDF page 4, Theorem 1.1, equations (1.1)-(1.2), Remark 1; arXiv 2301.06996v7}\par
{\footnotesize \textit{Scope.} External source text; selected passages actually read; local applications retain their own proof hypotheses}\par
{\footnotesize SHA-256 \texttt{418810dbbd1a1b8d} \quad reading copy: C:/WORKHOUSE/REPO/literature/inbox/JW\_\allowbreak{}2006/extracted/MONDAL\_\allowbreak{}2023\_\allowbreak{}GEOMETRIC.txt. The extracted-literature guide records its resolution \quad \textbf{workspace path under ALL THEORY/}}\par
\smallskip
\noindent\textbf{\texttt{guideLitOS1}} \hfill {\footnotesize outside this checkout}\par
{\small Axioms for Euclidean Green's functions}\par
{\footnotesize\ttfamily C:/\allowbreak{}WORKHOUSE/\allowbreak{}ALL THEORY/\allowbreak{}WORKHOUSE/\allowbreak{}literature/\allowbreak{}inbox/\allowbreak{}JW\_\allowbreak{}2006/\allowbreak{}extracted/\allowbreak{}OS\_\allowbreak{}1973.txt}\par
{\footnotesize \textit{Located.} PDF pages 1-3, introduction and axiom list; page 7 (printed 89), remark 2 after (3.1)}\par
{\footnotesize \textit{Scope.} External source text; selected passages actually read; local applications retain their own proof hypotheses}\par
{\footnotesize SHA-256 \texttt{12e78898f906a93f} \quad reading copy: C:/WORKHOUSE/REPO/literature/inbox/JW\_\allowbreak{}2006/extracted/OS\_\allowbreak{}1973.txt. The extracted-literature guide records its resolution \quad \textbf{workspace path under ALL THEORY/}}\par
\smallskip
\noindent\textbf{\texttt{guideLitOS2}} \hfill {\footnotesize outside this checkout}\par
{\small Axioms for Euclidean Green's functions II}\par
{\footnotesize\ttfamily C:/\allowbreak{}WORKHOUSE/\allowbreak{}ALL THEORY/\allowbreak{}WORKHOUSE/\allowbreak{}literature/\allowbreak{}inbox/\allowbreak{}JW\_\allowbreak{}2006/\allowbreak{}extracted/\allowbreak{}OS\_\allowbreak{}1975.txt}\par
{\footnotesize \textit{Located.} PDF pages 1-2, correction; page 7 (printed 287), IV.1, (4.1)-(4.2), Theorem E prime to R prime; page 11 (printed 291), (5.2)-(5.3) and totality}\par
{\footnotesize \textit{Scope.} External source text; selected passages actually read; local applications retain their own proof hypotheses}\par
{\footnotesize SHA-256 \texttt{93e83a4166240af0} \quad reading copy: C:/WORKHOUSE/REPO/literature/inbox/JW\_\allowbreak{}2006/extracted/OS\_\allowbreak{}1975.txt. The extracted-literature guide records its resolution \quad \textbf{workspace path under ALL THEORY/}}\par
\smallskip
\noindent\textbf{\texttt{guideLitPages}} \hfill {\footnotesize verified}\par
{\small Yang-Mills page-locator index}\par
{\footnotesize\ttfamily C:/\allowbreak{}WORKHOUSE/\allowbreak{}REPO/\allowbreak{}literature/\allowbreak{}yangmills/\allowbreak{}extraction/\allowbreak{}pages.jsonl}\par
{\footnotesize \textit{Located.} All 777 page records for 28 works}\par
{\footnotesize \textit{Scope.} Locator metadata only; no extracted full text}\par
{\footnotesize SHA-256 \texttt{af6f67963051dd93}}\par
\smallskip
\noindent\textbf{\texttt{guideLitReconstruction}} \hfill {\footnotesize verified}\par
{\small Reconstruction, observable completeness, and a quantitative localization interface}\par
{\footnotesize\ttfamily C:/\allowbreak{}WORKHOUSE/\allowbreak{}REPO/\allowbreak{}docs/\allowbreak{}derivations/\allowbreak{}yangmills-\allowbreak{}reconstruction.md}\par
{\footnotesize \textit{Located.} Sections 1-5, R1-R7; source statements on corrected OS hypotheses}\par
{\footnotesize \textit{Scope.} Local analytic derivation and finite calibration; remaining application premises explicit}\par
{\footnotesize SHA-256 \texttt{a180dddf14179838} \quad graph: \texttt{CITE:\allowbreak{}YM\_\allowbreak{}RECONSTRUCTION}}\par
\smallskip
\noindent\textbf{\texttt{guideLitSimon}} \hfill {\footnotesize outside this checkout}\par
{\small Some quantum operators with discrete spectrum but classically continuous spectrum}\par
{\footnotesize\ttfamily C:/\allowbreak{}WORKHOUSE/\allowbreak{}ALL THEORY/\allowbreak{}WORKHOUSE/\allowbreak{}literature/\allowbreak{}inbox/\allowbreak{}JW\_\allowbreak{}2006/\allowbreak{}extracted/\allowbreak{}SIMON\_\allowbreak{}1983\_\allowbreak{}DISCRETE.txt}\par
{\footnotesize \textit{Located.} PDF pages 1-3, introduction, equation (3), scalar mechanism; pages 9-10, Corollary 4 and matrix representation remark}\par
{\footnotesize \textit{Scope.} External source text; selected passages actually read; local applications retain their own proof hypotheses}\par
{\footnotesize SHA-256 \texttt{5112957632fe2e45} \quad reading copy: C:/WORKHOUSE/REPO/literature/inbox/JW\_\allowbreak{}2006/extracted/SIMON\_\allowbreak{}1983\_\allowbreak{}DISCRETE.txt. The extracted-literature guide records its resolution \quad \textbf{workspace path under ALL THEORY/}}\par
\smallskip
\noindent\textbf{\texttt{guideLitSimonBridge}} \hfill {\footnotesize verified}\par
{\small Yang-Mills flat directions: exact finite-dimensional confinement comparisons}\par
{\footnotesize\ttfamily C:/\allowbreak{}WORKHOUSE/\allowbreak{}REPO/\allowbreak{}docs/\allowbreak{}derivations/\allowbreak{}yangmills-\allowbreak{}simon-\allowbreak{}flat-\allowbreak{}directions.md}\par
{\footnotesize \textit{Located.} Scalar split, transverse oscillator, balanced aggregation, compact-resolvent and scaling arguments}\par
{\footnotesize \textit{Scope.} Local finite-model derivation implementing the published confinement mechanism}\par
{\footnotesize SHA-256 \texttt{850d8bf0ea4b6192} \quad graph: \texttt{CITE:\allowbreak{}YM\_\allowbreak{}FLAT\_\allowbreak{}DIRECTIONS}}\par
\smallskip
\noindent\textbf{\texttt{guideLitUeltschi}} \hfill {\footnotesize outside this checkout}\par
{\small Cluster expansions and correlation functions}\par
{\footnotesize\ttfamily C:/\allowbreak{}WORKHOUSE/\allowbreak{}ALL THEORY/\allowbreak{}WORKHOUSE/\allowbreak{}literature/\allowbreak{}inbox/\allowbreak{}WILSON\_\allowbreak{}MARKED\_\allowbreak{}20260908/\allowbreak{}UELTSCHI\_\allowbreak{}2004\_\allowbreak{}CLUSTER.txt}\par
{\footnotesize \textit{Located.} PDF pages 1-2, Theorem 1, equations (1)-(5); extract identifies arXiv math-ph/0304003v3}\par
{\footnotesize \textit{Scope.} External source text; selected passages actually read; local applications retain their own proof hypotheses}\par
{\footnotesize SHA-256 \texttt{fbe3771f7857839e} \quad reading copy: C:/WORKHOUSE/REPO/literature/inbox/WILSON\_\allowbreak{}MARKED\_\allowbreak{}20260908/UELTSCHI\_\allowbreak{}2004\_\allowbreak{}CLUSTER.txt. The extracted-literature guide records its resolution \quad \textbf{workspace path under ALL THEORY/}}\par
\smallskip
\noindent\textbf{\texttt{guideLitWilsonBlocks}} \hfill {\footnotesize verified}\par
{\small Actual Wilson transfer: blocks, vacuum cancellation and uniform estimates}\par
{\footnotesize\ttfamily C:/\allowbreak{}WORKHOUSE/\allowbreak{}REPO/\allowbreak{}docs/\allowbreak{}derivations/\allowbreak{}wilson-\allowbreak{}marked-\allowbreak{}transfer.md}\par
{\footnotesize \textit{Located.} WT1-WT5, especially WT2, WT3, WT4a-WT4b; continuation note precedes WT1}\par
{\footnotesize \textit{Scope.} Local discrete-time block application of the inspected cluster criterion; current shell continuation is separately cited}\par
{\footnotesize SHA-256 \texttt{e7096d99ced64935} \quad graph: \texttt{CITE:\allowbreak{}WILSON\_\allowbreak{}MARKED}}\par
\smallskip
\noindent\textbf{\texttt{guideLitYarotskyGS}} \hfill {\footnotesize outside this checkout}\par
{\small Ground states in relatively bounded quantum perturbations of classical lattice systems}\par
{\footnotesize\ttfamily C:/\allowbreak{}WORKHOUSE/\allowbreak{}ALL THEORY/\allowbreak{}WORKHOUSE/\allowbreak{}literature/\allowbreak{}inbox/\allowbreak{}WILSON\_\allowbreak{}MARKED\_\allowbreak{}20260908/\allowbreak{}YAROTSKY\_\allowbreak{}2004\_\allowbreak{}GS.txt}\par
{\footnotesize \textit{Located.} PDF pages 7-10, active-block expansion, Lemmas 1-4, equations (12)-(13), connected vacuum expansion}\par
{\footnotesize \textit{Scope.} External source text; selected passages actually read; local applications retain their own proof hypotheses}\par
{\footnotesize SHA-256 \texttt{c7f1c8180876d061} \quad reading copy: C:/WORKHOUSE/REPO/literature/inbox/WILSON\_\allowbreak{}MARKED\_\allowbreak{}20260908/YAROTSKY\_\allowbreak{}2004\_\allowbreak{}GS.txt. The extracted-literature guide records its resolution \quad \textbf{workspace path under ALL THEORY/}}\par
\smallskip
\noindent\textbf{\texttt{guideLitYarotskyQP}} \hfill {\footnotesize outside this checkout}\par
{\small Quasi-particles in weak perturbations of non-interacting quantum lattice systems}\par
{\footnotesize\ttfamily C:/\allowbreak{}WORKHOUSE/\allowbreak{}ALL THEORY/\allowbreak{}WORKHOUSE/\allowbreak{}literature/\allowbreak{}inbox/\allowbreak{}WILSON\_\allowbreak{}MARKED\_\allowbreak{}20260908/\allowbreak{}YAROTSKY\_\allowbreak{}2004\_\allowbreak{}QP.txt}\par
{\footnotesize \textit{Located.} PDF pages 2-5, model and Theorems 1-4; pages 6-10, equations (7)-(15) and start of section 3}\par
{\footnotesize \textit{Scope.} External source text; selected passages actually read; local applications retain their own proof hypotheses}\par
{\footnotesize SHA-256 \texttt{0089ae78e1c130ed} \quad reading copy: C:/WORKHOUSE/REPO/literature/inbox/WILSON\_\allowbreak{}MARKED\_\allowbreak{}20260908/YAROTSKY\_\allowbreak{}2004\_\allowbreak{}QP.txt. The extracted-literature guide records its resolution \quad \textbf{workspace path under ALL THEORY/}}\par
\smallskip
\noindent\textbf{\texttt{guideOperations}} \hfill {\footnotesize verified}\par
{\small Workspace operations for agents}\par
{\footnotesize\ttfamily C:/\allowbreak{}WORKHOUSE/\allowbreak{}REPO/\allowbreak{}docs/\allowbreak{}workspace\_\allowbreak{}operations.md}\par
{\footnotesize \textit{Located.} Select the working location; start with observed state; ownership and shared writers}\par
{\footnotesize \textit{Scope.} Operational source; no scientific theorem is inferred from the instructions}\par
{\footnotesize SHA-256 \texttt{4f0fe8ed751d3011}}\par
\smallskip
\noindent\textbf{\texttt{guideProtocol}} \hfill {\footnotesize verified}\par
{\small THEORY GRAPH protocol}\par
{\footnotesize\ttfamily C:/\allowbreak{}WORKHOUSE/\allowbreak{}REPO/\allowbreak{}docs/\allowbreak{}theory\_\allowbreak{}graph\_\allowbreak{}protocol.md}\par
{\footnotesize \textit{Located.} Retained snapshots; saved/live/fresh semantics; relationships; manual handoff}\par
{\footnotesize \textit{Scope.} Operational source; no scientific theorem is inferred from the instructions}\par
{\footnotesize SHA-256 \texttt{54cbbe206808fa30}}\par
\smallskip
\noindent\textbf{\texttt{guideTaskRecord}} \hfill {\footnotesize verified}\par
{\small Manual theory-graph task records}\par
{\footnotesize\ttfamily C:/\allowbreak{}WORKHOUSE/\allowbreak{}REPO/\allowbreak{}graph-\allowbreak{}tasks/\allowbreak{}README.md}\par
{\footnotesize \textit{Located.} Required fields; snapshot retention; unavailable automatic conformance checks}\par
{\footnotesize \textit{Scope.} Operational source; no scientific theorem is inferred from the instructions}\par
{\footnotesize SHA-256 \texttt{ebd725ff0732fd31}}\par
\smallskip
\noindent\textbf{\texttt{guideWorkspace}} \hfill {\footnotesize outside this checkout}\par
{\small Workspace identity and path roles}\par
{\footnotesize\ttfamily C:/\allowbreak{}WORKHOUSE/\allowbreak{}WORKSPACE.json}\par
{\footnotesize \textit{Located.} schema, canonical\_\allowbreak{}repository, pending\_\allowbreak{}research\_\allowbreak{}checkout and standalone\_\allowbreak{}research}\par
{\footnotesize \textit{Scope.} Operational source; no scientific theorem is inferred from the instructions}\par
{\footnotesize SHA-256 \texttt{eeae450bb2face85} \quad \textbf{workspace path under WORKSPACE.json/}}\par
\smallskip
\noindent\textbf{\texttt{hodgeexact}} \hfill {\footnotesize verified}\par
{\small Exact orbit-to-Hodge assembly of the charge-odd kernel}\par
{\footnotesize\ttfamily C:/\allowbreak{}WORKHOUSE/\allowbreak{}REPO/\allowbreak{}src/\allowbreak{}workhouse/\allowbreak{}kernel\_\allowbreak{}orbits.py}\par
{\footnotesize \textit{Located.} kernel\_\allowbreak{}orbits.py hodge\_\allowbreak{}form/hodge\_\allowbreak{}records, lines 743-772; constants.py raw orbit definitions, lines 447-488}\par
{\footnotesize \textit{Scope.} Exact raw-to-Hodge coefficient conversion and assembled fourth-order support at SU3}\par
{\footnotesize SHA-256 \texttt{ffc4f13643ff554f}}\par
\smallskip
\noindent\textbf{\texttt{localwick}} \hfill {\footnotesize verified}\par
{\small A rank-uniform continuation of the local SU(N) class spectrum}\par
{\footnotesize\ttfamily C:/\allowbreak{}WORKHOUSE/\allowbreak{}REPO/\allowbreak{}docs/\allowbreak{}derivations/\allowbreak{}local-\allowbreak{}class-\allowbreak{}wick-\allowbreak{}spectrum.md}\par
{\footnotesize \textit{Located.} Sections 1-6; recurrence, coefficients and signs}\par
{\footnotesize SHA-256 \texttt{5cde9ee4fe88bd47} \quad graph: \texttt{CITE:\allowbreak{}LOCAL\_\allowbreak{}CLASS\_\allowbreak{}WICK\_\allowbreak{}SPECTRUM}}\par
\smallskip
\noindent\textbf{\texttt{prior-carrier-bridge}} \hfill {\footnotesize verified}\par
{\small G18\_\allowbreak{}FIXED\_\allowbreak{}SPACING\_\allowbreak{}CARRIER\_\allowbreak{}BRIDGE\_\allowbreak{}INSERT.tex}\par
{\footnotesize\ttfamily C:/\allowbreak{}Users/\allowbreak{}Alex/\allowbreak{}Desktop/\allowbreak{}FINAL PAPERS/\allowbreak{}WORKHOUSE\_\allowbreak{}PUBLICATION\_\allowbreak{}EDITION\_\allowbreak{}FINAL\_\allowbreak{}20260830/\allowbreak{}research\_\allowbreak{}notes/\allowbreak{}G18\_\allowbreak{}FIXED\_\allowbreak{}SPACING\_\allowbreak{}CARRIER\_\allowbreak{}BRIDGE\_\allowbreak{}INSERT.tex}\par
{\footnotesize \textit{Located.} Fixed-band theorem at line 164; compact residue at line 370; input at line 406}\par
{\footnotesize \textit{Scope.} Established SU3 Hamiltonian patchwise carrier identification in the supplied August final edition}\par
{\footnotesize SHA-256 \texttt{87dc50be936d5b93}}\par
\smallskip
\noindent\textbf{\texttt{prior-carrier-sheet}} \hfill {\footnotesize verified}\par
{\small G18\_\allowbreak{}INTERNAL\_\allowbreak{}BBDAGGER\_\allowbreak{}SHEET\_\allowbreak{}CLOSURE\_\allowbreak{}20260830.tex}\par
{\footnotesize\ttfamily C:/\allowbreak{}Users/\allowbreak{}Alex/\allowbreak{}Desktop/\allowbreak{}FINAL PAPERS/\allowbreak{}WORKHOUSE\_\allowbreak{}PUBLICATION\_\allowbreak{}EDITION\_\allowbreak{}FINAL\_\allowbreak{}20260830/\allowbreak{}research\_\allowbreak{}notes/\allowbreak{}G18\_\allowbreak{}INTERNAL\_\allowbreak{}BBDAGGER\_\allowbreak{}SHEET\_\allowbreak{}CLOSURE\_\allowbreak{}20260830.tex}\par
{\footnotesize \textit{Located.} Patch theorem lines 309-396; no-dark corollary lines 398-453}\par
{\footnotesize \textit{Scope.} Established SU3 Hamiltonian patchwise carrier identification in the supplied August final edition}\par
{\footnotesize SHA-256 \texttt{8c871855b5bb415b}}\par
\smallskip
\noindent\textbf{\texttt{second-check-source-1}} \hfill {\footnotesize verified}\par
{\small constants.py}\par
{\footnotesize\ttfamily C:/\allowbreak{}WORKHOUSE/\allowbreak{}REPO/\allowbreak{}src/\allowbreak{}workhouse/\allowbreak{}constants.py}\par
{\footnotesize \textit{Located.} Source pinned by the second-pass exact mathematical check script}\par
{\footnotesize \textit{Scope.} Supporting exact check input; see evidence/microscopic\_\allowbreak{}second\_\allowbreak{}checks.py}\par
{\footnotesize SHA-256 \texttt{fd6580de321ec5ed}}\par
\smallskip
\noindent\textbf{\texttt{second-check-source-13}} \hfill {\footnotesize verified}\par
{\small coefficients\_\allowbreak{}order5.json}\par
{\footnotesize\ttfamily C:/\allowbreak{}WORKHOUSE/\allowbreak{}REPO/\allowbreak{}runs/\allowbreak{}local\_\allowbreak{}class\_\allowbreak{}wick\_\allowbreak{}2026-\allowbreak{}09-\allowbreak{}11/\allowbreak{}received/\allowbreak{}coefficients\_\allowbreak{}order5.json}\par
{\footnotesize \textit{Located.} Source pinned by the second-pass exact mathematical check script}\par
{\footnotesize \textit{Scope.} Supporting exact check input; see evidence/microscopic\_\allowbreak{}second\_\allowbreak{}checks.py}\par
{\footnotesize SHA-256 \texttt{82a0571580ec217d}}\par
\smallskip
\noindent\textbf{\texttt{second-check-source-3}} \hfill {\footnotesize verified}\par
{\small sixth\_\allowbreak{}order.py}\par
{\footnotesize\ttfamily C:/\allowbreak{}WORKHOUSE/\allowbreak{}REPO/\allowbreak{}src/\allowbreak{}workhouse/\allowbreak{}sixth\_\allowbreak{}order.py}\par
{\footnotesize \textit{Located.} Source pinned by the second-pass exact mathematical check script}\par
{\footnotesize \textit{Scope.} Supporting exact check input; see evidence/microscopic\_\allowbreak{}second\_\allowbreak{}checks.py}\par
{\footnotesize SHA-256 \texttt{76de82f6eeb16cd6}}\par
\smallskip
\noindent\textbf{\texttt{second-check-source-6}} \hfill {\footnotesize verified}\par
{\small electric\_\allowbreak{}shell.py}\par
{\footnotesize\ttfamily C:/\allowbreak{}WORKHOUSE/\allowbreak{}REPO/\allowbreak{}src/\allowbreak{}workhouse/\allowbreak{}invariants/\allowbreak{}electric\_\allowbreak{}shell.py}\par
{\footnotesize \textit{Located.} Source pinned by the second-pass exact mathematical check script}\par
{\footnotesize \textit{Scope.} Supporting exact check input; see evidence/microscopic\_\allowbreak{}second\_\allowbreak{}checks.py}\par
{\footnotesize SHA-256 \texttt{35b002a5b65bdcba}}\par
\smallskip
\noindent\textbf{\texttt{second-check-source-7}} \hfill {\footnotesize verified}\par
{\small local\_\allowbreak{}class\_\allowbreak{}wick.py}\par
{\footnotesize\ttfamily C:/\allowbreak{}WORKHOUSE/\allowbreak{}REPO/\allowbreak{}src/\allowbreak{}workhouse/\allowbreak{}local\_\allowbreak{}class\_\allowbreak{}wick.py}\par
{\footnotesize \textit{Located.} Source pinned by the second-pass exact mathematical check script}\par
{\footnotesize \textit{Scope.} Supporting exact check input; see evidence/microscopic\_\allowbreak{}second\_\allowbreak{}checks.py}\par
{\footnotesize SHA-256 \texttt{cbe081194bd68523}}\par
\smallskip
\noindent\textbf{\texttt{smith20260830}} \hfill {\footnotesize verified}\par
{\small Volume-uniform electric-shell isolation, exact shared-link hopping, and a homological carrier in strong-coupling SU(N) Hamiltonian lattice gauge theory}\par
{\footnotesize\ttfamily C:/\allowbreak{}Users/\allowbreak{}Alex/\allowbreak{}Desktop/\allowbreak{}FINAL PAPERS/\allowbreak{}WORKHOUSE\_\allowbreak{}PUBLICATION\_\allowbreak{}EDITION\_\allowbreak{}FINAL\_\allowbreak{}20260830/\allowbreak{}workhouse\_\allowbreak{}publication\_\allowbreak{}edition\_\allowbreak{}final\_\allowbreak{}20260830.tex}\par
{\footnotesize \textit{Located.} Final publication edition, 30 August 2026; manuscript and its listed inserts}\par
{\footnotesize SHA-256 \texttt{7d9022c8c4d49d95}}\par
\smallskip
\noindent\textbf{\texttt{swapodd}} \hfill {\footnotesize verified}\par
{\small ADR 0023: The swap-odd domino state derives leak = hop, and the vacuum route is not why}\par
{\footnotesize\ttfamily C:/\allowbreak{}WORKHOUSE/\allowbreak{}REPO/\allowbreak{}docs/\allowbreak{}decisions/\allowbreak{}0023-\allowbreak{}the-\allowbreak{}swap-\allowbreak{}odd-\allowbreak{}state-\allowbreak{}derives-\allowbreak{}the-\allowbreak{}c-\allowbreak{}even-\allowbreak{}identity.md}\par
{\footnotesize \textit{Located.} What was derived; Addendum disconnected fourth-order falsifier}\par
{\footnotesize \textit{Scope.} Charge-even equality through cubic order, disconnected 63/800 at quartic. Older vacuum T3 statement superseded by swapoddexact}\par
{\footnotesize SHA-256 \texttt{80bc2a5c760528fa} \quad graph: \texttt{ADR:0023}}\par
\smallskip
\noindent\textbf{\texttt{swapoddexact}} \hfill {\footnotesize verified}\par
{\small The swap-odd domino state: exact rotor towers and linked two-face vacuum}\par
{\footnotesize\ttfamily C:/\allowbreak{}WORKHOUSE/\allowbreak{}REPO/\allowbreak{}src/\allowbreak{}workhouse/\allowbreak{}invariants/\allowbreak{}swap\_\allowbreak{}odd.py}\par
{\footnotesize \textit{Located.} \_\allowbreak{}sector\_\allowbreak{}energies; fourth-order rotor check; \_\allowbreak{}connected\_\allowbreak{}pair\_\allowbreak{}vacuum\_\allowbreak{}su3; \_\allowbreak{}connected\_\allowbreak{}pair\_\allowbreak{}vacuum\_\allowbreak{}all\_\allowbreak{}rank; linked vacuum check at lines 441-508}\par
{\footnotesize \textit{Scope.} Exact SU3 omega4=-327/83776 plus Q(N) corroboration with min\_\allowbreak{}rank9; conn\_\allowbreak{}A remains open on (6,0) Haar family}\par
{\footnotesize SHA-256 \texttt{d8b06904f88e620e}}\par
\smallskip
\endgroup


\section{Statement-to-evidence crosswalk}
\label{app:crosswalk}

\CrosswalkRows{} numbered statements, each mapped to the source that
proves it, the labelled theorem inside that source, the graph entry
that records it, the scope of that evidence, and the file that checks
it.

Follow the source before treating a neighbouring graph node as the same
theorem. \texttt{RESULT} records an analytic result; \texttt{DERIV}
records a located source statement; a gap or a decision record is
research context and not evidence. None of the four is interchangeable,
and a crosswalk is the place that distinction is most often lost.

One consequence deserves stating separately, because it governs how the
\textit{Checked by} column should be read. A source ledger's
whole-statement formalization covers its entire named source statement;
a support entry covers only a recorded ingredient. That distinction
survives when several source statements are assembled into one theorem
of this document --- so a theorem here can be fully proved while no
single registered formalization covers it, and the absence of a
standalone registration below invalidates nothing.

Every graph entry in the table was resolved against the catalogue when
this document was built; \CrosswalkUnresolved{} did not resolve.

% Generated by paper/master/generate_sources.py.  Do not edit.
\noindent 36 numbered statements, each mapped to the source that proves it, the graph entry that records it, the scope of that evidence, and the file that checks it. Follow the source before treating a neighbouring graph node as the same theorem: \texttt{RESULT} is an analytic result, \texttt{DERIV} a located source statement, and a gap or decision is context rather than evidence. None of the types is interchangeable.\par\medskip
\begingroup\sloppy
\noindent\textbf{2. Second-order transport}\par
{\footnotesize \textit{Source.} \texttt{smith20260830}, \texttt{fourth20260909}}\par
{\footnotesize \textit{At.} \texttt{thm:weingarten}, \texttt{thm:allrank}, \texttt{lem:process-complete}, \texttt{thm:fact}, \texttt{thm:carrier}, \texttt{thm:bridge}}\par
{\footnotesize \textit{Graph.} \texttt{G24}}\par
{\footnotesize \textit{Scope.} Composite transport theorem. Lean checks the rational channel step and t3, not the full process exhaustion and physical Bloch identification}\par
{\footnotesize \textit{Checked by.} electric\_\allowbreak{}shell.py, rank\_\allowbreak{}law.py; 2 named Lean ingredients in JSON}\par
\smallskip
\noindent\textbf{3. The assembled fourth-order operator, \$SU(3)\$}\par
{\footnotesize \textit{Source.} \texttt{fourth20260909}, \texttt{hodgeexact}, \texttt{adr0024}}\par
{\footnotesize \textit{At.} \texttt{thm:hodge}, \texttt{thm:support}, \texttt{thm:c2}, \texttt{hodge\_\allowbreak{}form/hodge\_\allowbreak{}records}, \texttt{ADR 0024 findings 1-4}}\par
{\footnotesize \textit{Graph.} \texttt{RESULT:\allowbreak{}TIER\_\allowbreak{}COLLAPSE\_\allowbreak{}ACTUAL\_\allowbreak{}H4\_\allowbreak{}SUPPORT}}\par
{\footnotesize \textit{Scope.} Actual supported H4 and corrected coefficient dictionary. The arbitrary R-linear assertion is not used}\par
{\footnotesize \textit{Checked by.} hodge\_\allowbreak{}feshbach.py, kernel\_\allowbreak{}orbits.py}\par
\smallskip
\noindent\textbf{4. Coercivity of the corrected fourth-order polynomial}\par
{\footnotesize \textit{Source.} \texttt{contBudgets}}\par
{\footnotesize \textit{At.} \texttt{Q1}, \texttt{equations Q1-Q2}}\par
{\footnotesize \textit{Graph.} \texttt{RESULT:\allowbreak{}ASSEMBLED\_\allowbreak{}Q4\_\allowbreak{}COERCIVITY}}\par
{\footnotesize \textit{Scope.} Registered exact quadratic identity and extremal eigenvalues; not an all-orders Hamiltonian estimate}\par
{\footnotesize \textit{Checked by.} spectral\_\allowbreak{}budgets.py}\par
\smallskip
\noindent\textbf{5. Sharp fourth-order bounds on the physical carrier}\par
{\footnotesize \textit{Source.} \texttt{contBudgets}}\par
{\footnotesize \textit{At.} \texttt{Q1-Q2 plus paper mic:carriercoercivity}}\par
{\footnotesize \textit{Graph.} \texttt{RESULT:\allowbreak{}ASSEMBLED\_\allowbreak{}Q4\_\allowbreak{}COERCIVITY}}\par
{\footnotesize \textit{Scope.} Unregistered standalone paper corollary on the nonnegative momentum cone; listed graph statements provide its premise}\par
{\footnotesize \textit{Checked by.} spectral\_\allowbreak{}budgets.py}\par
\smallskip
\noindent\textbf{6. Only the non-Hodge generator excites the complement}\par
{\footnotesize \textit{Source.} \texttt{fourth20260909}, \texttt{g9combined}}\par
{\footnotesize \textit{At.} \texttt{prop:feshbach}, \texttt{G9 combined Section 3}}\par
{\footnotesize \textit{Graph.} \texttt{RESULT:\allowbreak{}HODGE\_\allowbreak{}FESHBACH\_\allowbreak{}SPLITTING}}\par
{\footnotesize \textit{Scope.} Actual H4 support and q>0 isolate R. Abstract Hodge Lean statements are ingredients}\par
{\footnotesize \textit{Checked by.} hodge\_\allowbreak{}feshbach.py; 4 named Lean ingredients in JSON}\par
\smallskip
\noindent\textbf{7. Hermitian sixth-order formula}\par
{\footnotesize \textit{Source.} \texttt{g9combined}}\par
{\footnotesize \textit{At.} \texttt{Section 2}, \texttt{equation F6}}\par
{\footnotesize \textit{Graph.} \texttt{RESULT:\allowbreak{}G9\_\allowbreak{}HERMITIAN\_\allowbreak{}FOLDED\_\allowbreak{}SUPPORT}}\par
{\footnotesize \textit{Scope.} Registered formal isolated-shell identity. Free-word and rational matrix checks do not identify Wilson matrix elements}\par
{\footnotesize \textit{Checked by.} sixth\_\allowbreak{}order.py}\par
\smallskip
\noindent\textbf{8. Sharp three-coordinate bound}\par
{\footnotesize \textit{Source.} \texttt{anisotropy}}\par
{\footnotesize \textit{At.} \texttt{TG1 constrained-extremum proof}}\par
{\footnotesize \textit{Graph.} \texttt{RESULT:\allowbreak{}ANISOTROPY\_\allowbreak{}THREE\_\allowbreak{}COORDINATE\_\allowbreak{}SHARP\_\allowbreak{}BOUND}}\par
{\footnotesize \textit{Scope.} Analytic sharp extremum. Lean variance identity and nonnegativity do not prove this maximum}\par
{\footnotesize \textit{Checked by.} theory\_\allowbreak{}current\_\allowbreak{}bridges.py; 2 named Lean ingredients in JSON}\par
\smallskip
\noindent\textbf{9. A surviving sixth-order rational shape}\par
{\footnotesize \textit{Source.} \texttt{g9combined}}\par
{\footnotesize \textit{At.} \texttt{Section 5}, \texttt{N6-R6 and complex-torus witness}}\par
{\footnotesize \textit{Graph.} \texttt{RESULT:\allowbreak{}G9\_\allowbreak{}LOCAL\_\allowbreak{}SIXTH\_\allowbreak{}NONCANCELLATION}}\par
{\footnotesize \textit{Scope.} Locality and whole-zone Laurent divisibility force noncancellation. Symmetric F is an extra premise only for the symmetric remainder}\par
{\footnotesize \textit{Checked by.} sixth\_\allowbreak{}order.py}\par
\smallskip
\noindent\textbf{10. Odd orders are determinant families}\par
{\footnotesize \textit{Source.} \texttt{adr0047}}\par
{\footnotesize \textit{At.} \texttt{Sections 2-3}, \texttt{centre parity and determinant path}}\par
{\footnotesize \textit{Graph.} \texttt{RESULT:\allowbreak{}ODD\_\allowbreak{}ORDER\_\allowbreak{}CENTRE\_\allowbreak{}PARITY}}\par
{\footnotesize \textit{Scope.} Analytic all-order theorem on the stated clusters and periodic extents. Finite-rank character/word controls have bounded scope}\par
{\footnotesize \textit{Checked by.} odd\_\allowbreak{}order.py}\par
\smallskip
\noindent\textbf{11. Finite polynomial shell resolution}\par
{\footnotesize \textit{Source.} \texttt{localwick}}\par
{\footnotesize \textit{At.} \texttt{Sections 1-2}, \texttt{Wick conjugation and shell projectors}}\par
{\footnotesize \textit{Graph.} \texttt{DERIV:\allowbreak{}LOCAL\_\allowbreak{}CLASS\_\allowbreak{}WICK:\allowbreak{}RECURRENCE}}\par
{\footnotesize \textit{Scope.} Source lemma within a registered construction; no distinct RESULT/DERIV for the paper's shell-resolution proposition}\par
{\footnotesize \textit{Checked by.} local\_\allowbreak{}class\_\allowbreak{}wick.py}\par
\smallskip
\noindent\textbf{12. Rank-uniform Wick recurrence}\par
{\footnotesize \textit{Source.} \texttt{localwick}}\par
{\footnotesize \textit{At.} \texttt{Section 3 all-order constructive theorem}}\par
{\footnotesize \textit{Graph.} \texttt{RESULT:\allowbreak{}LOCAL\_\allowbreak{}CLASS\_\allowbreak{}WICK\_\allowbreak{}RECURRENCE}}\par
{\footnotesize \textit{Scope.} Analytic formal all-rank recurrence. Order-five residuals and trace controls do not establish a compact-group remainder}\par
{\footnotesize \textit{Checked by.} local\_\allowbreak{}class\_\allowbreak{}wick.py}\par
\smallskip
\noindent\textbf{13. Signs of the first five corrections}\par
{\footnotesize \textit{Source.} \texttt{localwick}}\par
{\footnotesize \textit{At.} \texttt{Section 4.2 sign corollary}}\par
{\footnotesize \textit{Graph.} \texttt{RESULT:\allowbreak{}LOCAL\_\allowbreak{}CLASS\_\allowbreak{}WICK\_\allowbreak{}SIGNS}}\par
{\footnotesize \textit{Scope.} Exact symbolic signs at admissible ranks, restricted to the displayed five coefficients}\par
{\footnotesize \textit{Checked by.} local\_\allowbreak{}class\_\allowbreak{}wick.py}\par
\smallskip
\noindent\textbf{14. Exact determinant-family quotients}\par
{\footnotesize \textit{Source.} \texttt{fourthMicroEngine}, \texttt{fourthMicroEpsilon}, \texttt{fourthMicroEpsilonTests}}\par
{\footnotesize \textit{At.} \texttt{Engine lines 810-967}, \texttt{paper Schur-Weyl completeness argument}}\par
{\footnotesize \textit{Graph.} \texttt{ADR:0021}}\par
{\footnotesize \textit{Scope.} Unregistered standalone paper formulation of pinned tensor algebra. ADR/G IDs locate the integration campaign, not a separate projector theorem}\par
{\footnotesize \textit{Checked by.} DATA\_\allowbreak{}SU3\_\allowbreak{}Exact\_\allowbreak{}MarkedCluster\_\allowbreak{}m4\_\allowbreak{}Colab.py, haar\_\allowbreak{}epsilon.py, test\_\allowbreak{}haar\_\allowbreak{}epsilon.py}\par
\smallskip
\noindent\textbf{15. Energy-preserving path reduction}\par
{\footnotesize \textit{Source.} \texttt{fourthMicroPaths}, \texttt{fourthMicroLoopcalc}, \texttt{fourthMicroPathRecord}}\par
{\footnotesize \textit{At.} \texttt{ADR 0030 path lemma}, \texttt{loopcalc lines 54-64}, \texttt{291-337}, \texttt{593-635}}\par
{\footnotesize \textit{Graph.} \texttt{ADR:0030}}\par
{\footnotesize \textit{Scope.} Paper operator formulation of the registered path argument. Product dependence, electric weight and compatible spectral projections are explicit; no standalone RESULT/DERIV}\par
{\footnotesize \textit{Checked by.} path\_\allowbreak{}reduction.py}\par
\smallskip
\noindent\textbf{16. The cellular carrier and its excursions}\par
{\footnotesize \textit{Source.} \texttt{fourthMicroCellular}}\par
{\footnotesize \textit{At.} \texttt{Sections 2.2-2.4}, \texttt{Theorems 1-3}}\par
{\footnotesize \textit{Graph.} \texttt{RESULT:\allowbreak{}UNIVERSAL\_\allowbreak{}CELLULAR\_\allowbreak{}HODGE\_\allowbreak{}SPECTRUM}}\par
{\footnotesize \textit{Scope.} Composite cellular theorem. Whole Lean coverage exists for the named rank-one excursion substatement only, not the complete incidence/perimeter theorem}\par
{\footnotesize \textit{Checked by.} universal\_\allowbreak{}cellular\_\allowbreak{}hodge.py; 2 named Lean ingredients in JSON}\par
\smallskip
\noindent\textbf{17. The complete tetrahedral face commutant}\par
{\footnotesize \textit{Source.} \texttt{fourthMicroCellular}}\par
{\footnotesize \textit{At.} \texttt{Sections 3.1-3.2}, \texttt{Theorems 4-5}}\par
{\footnotesize \textit{Graph.} \texttt{RESULT:\allowbreak{}TETRAHEDRAL\_\allowbreak{}S4\_\allowbreak{}COMMUTANT\_\allowbreak{}RESOLUTION}}\par
{\footnotesize \textit{Scope.} Analytic full S4 commutant. Lean duality/scalar trace helpers are scoped; finite checks solve all 24 permutations}\par
{\footnotesize \textit{Checked by.} universal\_\allowbreak{}cellular\_\allowbreak{}hodge.py; 2 named Lean ingredients in JSON}\par
\smallskip
\noindent\textbf{18. First connected isotropic cap dispersion}\par
{\footnotesize \textit{Source.} \texttt{fourthCapMaster}, \texttt{fourthCapRun}, \texttt{fourthCapChecks}}\par
{\footnotesize \textit{At.} \texttt{Master Section 9.3}, \texttt{cold-run What it produced and What it changes}}\par
{\footnotesize \textit{Graph.} \texttt{CHK:isotropic-pentagonal-cap-band-v4-3-9-3:h-4-side--a---a--exactly}}\par
{\footnotesize \textit{Scope.} Legacy separate-geometry theorem with native exact checks; no distinct RESULT/DERIV identified. Local-open/periodic L=5 support and continuous-zone bandwidth remain distinct}\par
{\footnotesize \textit{Checked by.} pentagonal.py}\par
\smallskip
\noindent\textbf{19. Complete physical Wilson band and onto literal sources}\par
{\footnotesize \textit{Source.} \texttt{contG18}}\par
{\footnotesize \textit{At.} \texttt{Sections 1-8}, \texttt{operator limit}, \texttt{tagged frame}, \texttt{rotation and physical history representation}}\par
{\footnotesize \textit{Graph.} \texttt{RESULT:\allowbreak{}WILSON\_\allowbreak{}INFINITE\_\allowbreak{}PHYSICAL\_\allowbreak{}BAND}}\par
{\footnotesize \textit{Scope.} Full infinite-dimensional theorem is analytic. Registered finite Gram, source, projection and activity checks are supporting controls}\par
{\footnotesize \textit{Checked by.} wilson\_\allowbreak{}physical\_\allowbreak{}band.py}\par
\smallskip
\noindent\textbf{20. SC17 thermodynamic physical-time construction}\par
{\footnotesize \textit{Source.} \texttt{contSC17}, \texttt{contSC17Time}, \texttt{fourthAnalyticThermo}}\par
{\footnotesize \textit{At.} \texttt{Maximized bootstrap}, \texttt{thermodynamic T12-T13}, \texttt{score limit and IF4-IF9}, \texttt{physical-time P1-P20}}\par
{\footnotesize \textit{Graph.} \texttt{DERIV:\allowbreak{}WILSON\_\allowbreak{}SC17\_\allowbreak{}SPATIAL\_\allowbreak{}CLOSURE:\allowbreak{}OPTIMIZED\_\allowbreak{}INTERVAL}}\par
{\footnotesize \textit{Scope.} Composite analytic endpoint theorem. Lean closure/gap-limit helpers do not identify the actual physical process or its semigroup}\par
{\footnotesize \textit{Checked by.} analytic source argument; 3 named Lean ingredients in JSON}\par
\smallskip
\noindent\textbf{21. Selected inverse from the complete residual}\par
{\footnotesize \textit{Source.} \texttt{contCurrents}}\par
{\footnotesize \textit{At.} \texttt{SCB3 and paper weak-form proof}}\par
{\footnotesize \textit{Graph.} \texttt{DERIV:\allowbreak{}SOURCE\_\allowbreak{}CURRENT\_\allowbreak{}BRIDGES:\allowbreak{}SCB3}}\par
{\footnotesize \textit{Scope.} Paper weak-form extension of a registered selected-inverse implication. The bounded-inverse ingredient is scoped support}\par
{\footnotesize \textit{Checked by.} theory\_\allowbreak{}current\_\allowbreak{}bridges.py; 1 named Lean ingredient in JSON}\par
\smallskip
\noindent\textbf{22. Approximate sources at one late physical time}\par
{\footnotesize \textit{Source.} \texttt{contMoving}}\par
{\footnotesize \textit{At.} \texttt{MT1-MT2}}\par
{\footnotesize \textit{Graph.} \texttt{RESULT:\allowbreak{}MOVING\_\allowbreak{}TIME\_\allowbreak{}SPECTRAL\_\allowbreak{}GAP}}\par
{\footnotesize \textit{Scope.} Analytic projection/vague-limit implication; finite positive-transfer and counterexample controls do not certify the limiting theorem}\par
{\footnotesize \textit{Checked by.} moving\_\allowbreak{}time\_\allowbreak{}gap.py}\par
\smallskip
\noindent\textbf{23. Hamiltonian reconstruction from compatible positive-time kernels}\par
{\footnotesize \textit{Source.} \texttt{contKernels}}\par
{\footnotesize \textit{At.} \texttt{K1 and K5}}\par
{\footnotesize \textit{Graph.} \texttt{RESULT:\allowbreak{}OS\_\allowbreak{}KERNEL\_\allowbreak{}MOVING\_\allowbreak{}TIME\_\allowbreak{}GAP}}\par
{\footnotesize \textit{Scope.} Analytic null quotient, positive-time smoothing and Hamiltonian reconstruction; finite Gram examples are scoped}\par
{\footnotesize \textit{Checked by.} os\_\allowbreak{}kernel\_\allowbreak{}gap.py}\par
\smallskip
\noindent\textbf{24. Common moving-time rate on the full history space}\par
{\footnotesize \textit{Source.} \texttt{contKernels}}\par
{\footnotesize \textit{At.} \texttt{K2 common-rate conclusion}}\par
{\footnotesize \textit{Graph.} \texttt{RESULT:\allowbreak{}OS\_\allowbreak{}KERNEL\_\allowbreak{}MOVING\_\allowbreak{}TIME\_\allowbreak{}GAP}}\par
{\footnotesize \textit{Scope.} Paper corollary of a registered source theorem, without a distinct RESULT ID. One common positive rate on the spanning family is essential}\par
{\footnotesize \textit{Checked by.} os\_\allowbreak{}kernel\_\allowbreak{}gap.py}\par
\smallskip
\noindent\textbf{25. A separated-time correlator retains finite-energy weight}\par
{\footnotesize \textit{Source.} \texttt{contKernels}}\par
{\footnotesize \textit{At.} \texttt{K3 finite-energy weight}}\par
{\footnotesize \textit{Graph.} \texttt{RESULT:\allowbreak{}OS\_\allowbreak{}KERNEL\_\allowbreak{}MOVING\_\allowbreak{}TIME\_\allowbreak{}GAP}}\par
{\footnotesize \textit{Scope.} Analytic positive-measure tail split; exact finite measures check the coefficient and denominator}\par
{\footnotesize \textit{Checked by.} os\_\allowbreak{}kernel\_\allowbreak{}gap.py}\par
\smallskip
\noindent\textbf{26. Uniform flat-background spatial source bound}\par
{\footnotesize \textit{Source.} \texttt{fourthHolonomy}, \texttt{fourthHolonomyLocal}, \texttt{fourthHolonomyControls}}\par
{\footnotesize \textit{At.} \texttt{Sections 1-3}, \texttt{equations 1-13}, \texttt{Local Sections 4-5}}\par
{\footnotesize \textit{Graph.} \texttt{RESULT:\allowbreak{}WILSON\_\allowbreak{}FLAT\_\allowbreak{}BACKGROUND\_\allowbreak{}FAST\_\allowbreak{}SOURCES}}\par
{\footnotesize \textit{Scope.} Analytic all-volume/all-flat-background source form and ranks. Retained exact Q(i) runs cover their specified finite backgrounds}\par
{\footnotesize \textit{Checked by.} analytic source argument}\par
\smallskip
\noindent\textbf{27. The physical source rank changes at central winding}\par
{\footnotesize \textit{Source.} \texttt{fourthHolonomy}, \texttt{fourthHolonomyControls}}\par
{\footnotesize \textit{At.} \texttt{Section 5}, \texttt{equations 17-18}}\par
{\footnotesize \textit{Graph.} \texttt{RESULT:\allowbreak{}WILSON\_\allowbreak{}PHYSICAL\_\allowbreak{}SOURCE\_\allowbreak{}RANK\_\allowbreak{}REPAIR}}\par
{\footnotesize \textit{Scope.} Analytic norm-one rank jump on the full SU2 family. Finite charged-vector and rank witnesses support it}\par
{\footnotesize \textit{Checked by.} analytic source argument}\par
\smallskip
\noindent\textbf{28. Uniform redundant lift and a nonflat submersion neighborhood}\par
{\footnotesize \textit{Source.} \texttt{fourthHolonomy}, \texttt{fourthHolonomyLocal}, \texttt{fourthHolonomyControls}}\par
{\footnotesize \textit{At.} \texttt{Section 4}, \texttt{equations 14-16 and 15b}}\par
{\footnotesize \textit{Graph.} \texttt{RESULT:\allowbreak{}WILSON\_\allowbreak{}FLAT\_\allowbreak{}BACKGROUND\_\allowbreak{}FAST\_\allowbreak{}SOURCES}}\par
{\footnotesize \textit{Scope.} Analytic right inverse, smooth redundant range and nonflat submersion. Finite boundary controls give no neighborhood probability}\par
{\footnotesize \textit{Checked by.} analytic source argument}\par
\smallskip
\noindent\textbf{29. Complete reference-graph Schur identity}\par
{\footnotesize \textit{Source.} \texttt{contScale}}\par
{\footnotesize \textit{At.} \texttt{Section 2}, \texttt{SP1-SP6}}\par
{\footnotesize \textit{Graph.} \texttt{DERIV:\allowbreak{}WILSON\_\allowbreak{}SPATIAL\_\allowbreak{}SCHUR\_\allowbreak{}EXCESS:\allowbreak{}SP1\_\allowbreak{}SP6}}\par
{\footnotesize \textit{Scope.} Analytic closed-form identity. Lean bounded square completion and minimizer are scoped, not the complete source/domain theorem}\par
{\footnotesize \textit{Checked by.} wilson\_\allowbreak{}spatial.py; 5 named Lean ingredients in JSON}\par
\smallskip
\noindent\textbf{30. Complete second-order excess with a cubic remainder}\par
{\footnotesize \textit{Source.} \texttt{contScale}}\par
{\footnotesize \textit{At.} \texttt{Section 3}, \texttt{SP7-SP10}}\par
{\footnotesize \textit{Graph.} \texttt{DERIV:\allowbreak{}WILSON\_\allowbreak{}SPATIAL\_\allowbreak{}SCHUR\_\allowbreak{}EXCESS:\allowbreak{}SP7\_\allowbreak{}SP10}}\par
{\footnotesize \textit{Scope.} Analytic cubic remainder from the stated uniform relative-form jet hypotheses; noncommuting finite controls check cross terms/constants}\par
{\footnotesize \textit{Checked by.} wilson\_\allowbreak{}spatial.py; 5 named Lean ingredients in JSON}\par
\smallskip
\noindent\textbf{31. Local response without a boundary-area factor}\par
{\footnotesize \textit{Source.} \texttt{fourthAnalyticThermo}}\par
{\footnotesize \textit{At.} \texttt{Part I}, \texttt{T3-T4 (T4-T11)}}\par
{\footnotesize \textit{Graph.} \texttt{DERIV:\allowbreak{}WILSON\_\allowbreak{}SC17\_\allowbreak{}THERMODYNAMIC\_\allowbreak{}LIMIT:\allowbreak{}T4\_\allowbreak{}T10}}\par
{\footnotesize \textit{Scope.} Analytic actual normalized cut response and covariance proof without boundary-area loss; no whole-statement Lean registration}\par
{\footnotesize \textit{Checked by.} analytic source argument}\par
\smallskip
\noindent\textbf{32. An explicit actual plaquette spectral interval}\par
{\footnotesize \textit{Source.} \texttt{fourthAnalyticThermo}}\par
{\footnotesize \textit{At.} \texttt{Part II Section 6}, \texttt{IF10-IF13}}\par
{\footnotesize \textit{Graph.} \texttt{DERIV:\allowbreak{}WILSON\_\allowbreak{}SC17\_\allowbreak{}THERMODYNAMIC\_\allowbreak{}LIMIT:\allowbreak{}IF10\_\allowbreak{}IF13}}\par
{\footnotesize \textit{Scope.} Lean proves variance domination and the supplied spectral-measure interval implication; actual Wilson density/Haar/observable identification remains analytic}\par
{\footnotesize \textit{Checked by.} analytic source argument; 4 named Lean ingredients in JSON}\par
\smallskip
\noindent\textbf{33. A common spatially weighted Taylor majorant}\par
{\footnotesize \textit{Source.} \texttt{fourthAnalyticTemporal}}\par
{\footnotesize \textit{At.} \texttt{Section 5}, \texttt{S16-S17}}\par
{\footnotesize \textit{Graph.} \texttt{DERIV:\allowbreak{}WILSON\_\allowbreak{}MARKED\_\allowbreak{}SHELL\_\allowbreak{}TRANSPORT:\allowbreak{}S16\_\allowbreak{}S17}}\par
{\footnotesize \textit{Scope.} Analytic common weighted majorant with actual holomorphic shell/frame; the native range-counting control is a scalar ingredient}\par
{\footnotesize \textit{Checked by.} wilson\_\allowbreak{}shell.py}\par
\smallskip
\noindent\textbf{34. Summed temporal matching and relative carrier transport}\par
{\footnotesize \textit{Source.} \texttt{fourthAnalyticTemporal}}\par
{\footnotesize \textit{At.} \texttt{Sections 6-7}, \texttt{S18-S22}}\par
{\footnotesize \textit{Graph.} \texttt{DERIV:\allowbreak{}WILSON\_\allowbreak{}MARKED\_\allowbreak{}SHELL\_\allowbreak{}TRANSPORT:\allowbreak{}S18\_\allowbreak{}S20}}\par
{\footnotesize \textit{Scope.} Analytic summed matching and specified-source transport; algebra controls do not identify overlap for arbitrary observables}\par
{\footnotesize \textit{Checked by.} wilson\_\allowbreak{}marked.py, wilson\_\allowbreak{}shell.py}\par
\smallskip
\noindent\textbf{35. Lyapunov conversion to a global Poincare bound}\par
{\footnotesize \textit{Source.} \texttt{anc-functional}}\par
{\footnotesize \textit{At.} \texttt{F8-F9}, \texttt{completed square and patching}}\par
{\footnotesize \textit{Graph.} \texttt{DERIV:\allowbreak{}ARCHIVE\_\allowbreak{}FUNCTIONAL:\allowbreak{}F8}}\par
{\footnotesize \textit{Scope.} Analytic drift-to-energy/local-to-global implication; the native exact-square check is only its algebraic ingredient}\par
{\footnotesize \textit{Checked by.} archive\_\allowbreak{}derivations.py}\par
\smallskip
\noindent\textbf{36. Accumulated physical gap budget}\par
{\footnotesize \textit{Source.} \texttt{contScale}}\par
{\footnotesize \textit{At.} \texttt{Section 6}, \texttt{SP20-SP22 plus paper geometric-series specialization}}\par
{\footnotesize \textit{Graph.} \texttt{DERIV:\allowbreak{}WILSON\_\allowbreak{}SPATIAL\_\allowbreak{}SCHUR\_\allowbreak{}EXCESS:\allowbreak{}SP20\_\allowbreak{}SP24}}\par
{\footnotesize \textit{Scope.} Paper specialization of registered scale recursion. Physical one-step comparisons are hypotheses; finite controls check the telescoping arithmetic}\par
{\footnotesize \textit{Checked by.} wilson\_\allowbreak{}spatial.py}\par
\smallskip
\noindent\textbf{--. }\par
{\footnotesize \textit{Graph.} \texttt{G17}}\par
{\footnotesize \textit{Scope.} Located proof, graph entry and evidence scope\textbackslash{}\textbackslash{} \textbackslash{}midrule\textbackslash{}endfirsthead \textbackslash{}toprule Paper statement \& Located proof, graph entry and evidence scope\textbackslash{}\textbackslash{} \textbackslash{}midrule\textbackslash{}endhead \textbackslash{}bottomrule\textbackslash{}endfoot 1. Proposition\textasciitilde{}mic:electricwindow The free electric window \& Source: smith20260830. thm:retained-shell; eq:first-window; eq:next-casimir; eq:zero-nality-casimir Legacy analytic Casimir/short-cycle theorem; G17 is context, and the native shell checks are scoped evidence}\par
{\footnotesize \textit{Checked by.} electric\_\allowbreak{}shell.py}\par
\smallskip
\endgroup

